\chapter{Section graded ring infrastructure: tensor powers and graded sections}
\label{chap:Picard_SectionGradedRing}

\providecommand{\PshMod}{\mathrm{PresheafOfModules}}
\providecommand{\ShMod}{\mathrm{SheafOfModules}}
\providecommand{\Modules}{\mathrm{Modules}}

The Hilbert-polynomial route of the parent Quot-scheme construction represents the
degree-\(m\) graded component of a fibre by the global sections
\(\Gamma(X_s, \mathcal{F}_s \otimes L_s^{\otimes m})\) of a twist of a coherent
sheaf by a tensor power of a line bundle, and it organises these components into
a graded ring \(R(X_s, L_s) = \bigoplus_{m \ge 0} \Gamma(X_s, L_s^{\otimes m})\)
(\cref{def:sectionGradedRing}) and a graded module
\(M(X_s, \mathcal{F}_s, L_s) = \bigoplus_{m \ge 0}
\Gamma(X_s, \mathcal{F}_s \otimes L_s^{\otimes m})\)
(\cref{def:sectionGradedModule}). Three pieces of that picture have no direct
Mathlib counterpart and must be built here, each in its own section:

\begin{enumerate}
  \item \textbf{Tensor powers of a sheaf of modules}
    (\cref{sec:sgr_tensor_powers}). Mathlib equips the category
    \(\PshMod\) of \emph{presheaves} of modules over a presheaf of commutative
    rings with a symmetric monoidal structure
    (\cref{lem:presheafModule_monoidal_mathlib}) and provides the
    sheafification functor \(\PshMod \to \ShMod\)
    (\cref{lem:presheafModule_sheafification_mathlib}); but the category
    \(\ShMod\) of \emph{sheaves} of modules carries \emph{no} monoidal
    structure. The tensor product of sheaves of modules
    \(\mathcal{F} \otimes_{\mathcal{O}_X} \mathcal{G}\) must therefore be built
    as the sheafification of the pointwise tensor presheaf, and the iterated
    tensor power \(\mathcal{L}^{\otimes m}\) recursively from it.
  \item \textbf{Lax-monoidal global sections} (\cref{sec:sgr_lax_sections}).
    The global-sections functor \(\Gamma\) is only \emph{lax} monoidal: from the
    pointwise-strict monoidality of the presheaf tensor product and the
    sheafification unit one obtains a natural multiplication
    \(\Gamma(\mathcal{F}) \otimes_{\Gamma(\mathcal{O}_X)} \Gamma(\mathcal{G})
    \to \Gamma(\mathcal{F} \otimes_{\mathcal{O}_X} \mathcal{G})\), together with
    the associativity and unit coherence those products must satisfy.
  \item \textbf{Graded-ring / graded-module assembly}
    (\cref{sec:sgr_graded_assembly}). The section multiplications and the
    tensor-power comparison isomorphisms assemble, through Mathlib's
    external-direct-sum graded API (\cref{lem:directSum_gcommSemiring_mathlib},
    \cref{lem:directSum_gmodule_mathlib}), into the graded \emph{ring} structure
    on \(\bigoplus_{m \ge 0} \Gamma(X_s, L_s^{\otimes m})\) and the graded
    \emph{module} structure on
    \(\bigoplus_{m \ge 0} \Gamma(X_s, \mathcal{F}_s \otimes L_s^{\otimes m})\).
\end{enumerate}

\section{Tensor powers of a sheaf of modules}
\label{sec:sgr_tensor_powers}

We work over a fixed scheme \(X\) (in the application \(X = X_s\), a fibre over a
field \(\kappa(s)\)). Its category of sheaves of modules is
\(X.\Modules = \ShMod(\mathcal{O}_X)\), with global-sections notation
\(\Gamma(X, \mathcal{M}) = \mathcal{M}(X)\) for the value of the underlying
presheaf on the top open. The structure sheaf \(\mathcal{O}_X\) is a sheaf of
commutative rings, so its underlying presheaf factors through the category of
commutative rings; this is the hypothesis under which Mathlib's presheaf-level
monoidal structure applies.

\begin{lemma}[Monoidal structure on presheaves of modules]
  \label{lem:presheafModule_monoidal_mathlib}
  \lean{PresheafOfModules.monoidalCategory}
  \mathlibok
  \textit{Provided by Mathlib
  (\texttt{Mathlib.Algebra.Category.ModuleCat.Presheaf.Monoidal}).}
  Let \(R : \mathcal{C}^{\mathrm{op}} \to \mathrm{CommRingCat}\) be a presheaf of
  commutative rings. Then the category \(\PshMod(R)\) of presheaves of modules
  over \(R\) carries a symmetric monoidal structure whose tensor product is
  computed objectwise: for presheaves of modules \(\mathcal{M}_1, \mathcal{M}_2\)
  and an object \(U\),
  \[
    (\mathcal{M}_1 \otimes \mathcal{M}_2)(U)
      \;=\; \mathcal{M}_1(U) \otimes_{R(U)} \mathcal{M}_2(U),
  \]
  with restriction maps the tensor products of the restriction maps. The unit is
  the presheaf \(R\) viewed as a module over itself, and the associator,
  unitors, and braiding are induced objectwise from those of
  \(\mathrm{ModuleCat}\). In particular evaluation at a fixed object \(U\) is a
  \emph{strict} monoidal functor
  \(\PshMod(R) \to \mathrm{ModuleCat}(R(U))\).
\end{lemma}

\begin{lemma}[Sheafification of presheaves of modules]
  \label{lem:presheafModule_sheafification_mathlib}
  \lean{PresheafOfModules.sheafification}
  \mathlibok
  \textit{Provided by Mathlib
  (\texttt{Mathlib.Algebra.Category.ModuleCat.Presheaf.Sheafification}).}
  Let \(R\) be a sheaf of (commutative) rings on a site, with underlying
  presheaf \(R_0\). There is a sheafification functor
  \[
    (-)^{\#} \;:\; \PshMod(R_0) \longrightarrow \ShMod(R),
  \]
  left adjoint to the forgetful functor, together with the unit
  \(\eta_{\mathcal{M}_0} : \mathcal{M}_0 \to (\mathcal{M}_0)^{\#}\) of the
  adjunction. On underlying sheaves of abelian groups it agrees with the
  ordinary sheafification, and it sends a presheaf of \(R_0\)-modules to the
  associated sheaf of \(R\)-modules.
\end{lemma}

\begin{definition}\leanok
[Module sheafification for a scheme]
  \label{def:schemeModuleSheafification}
  \lean{AlgebraicGeometry.Scheme.Modules.sheafification}
  \uses{lem:presheafModule_sheafification_mathlib}
  For a scheme \(X\) the \emph{module sheafification functor}
  \[
    (-)^{\#} \;:\; \PshMod(\mathcal{O}_X) \longrightarrow X.\Modules
  \]
  is the specialisation of Mathlib's presheaf-of-modules sheafification
  (\cref{lem:presheafModule_sheafification_mathlib}) to the identity morphism of
  the underlying presheaf of rings \(\mathcal{O}_X\) (which is locally bijective).
  It is the left adjoint of the forgetful functor
  \(X.\Modules \to \PshMod(\mathcal{O}_X)\); the tensor product of sheaves of
  modules (\cref{def:sheafTensorObj}) is by definition the value of this functor
  on the presheaf tensor product, and the unitor and braiding isomorphisms
  (\cref{def:tensorObjUnitIso}, \cref{def:tensorBraiding}) are obtained by
  applying it to the corresponding presheaf-level isomorphisms. As a reflective
  localization it inverts exactly the morphisms whose underlying abelian-presheaf
  morphism is a local isomorphism (\cref{lem:isIso_sheafification_map_iff}).
\end{definition}

\begin{lemma}[The structure sheaf as the unit module]
  \label{lem:moduleUnit_mathlib}
  \lean{SheafOfModules.unit}
  \mathlibok
  \textit{Provided by Mathlib
  (\texttt{Mathlib.Algebra.Category.ModuleCat.Sheaf}).}
  For a sheaf of rings \(R\) (in the application \(R = \mathcal{O}_X\)) the
  structure sheaf, viewed as a module over itself, is a distinguished object
  \(\mathbf{1} = \mathcal{O}_X \in \ShMod(\mathcal{O}_X)\). It is the value at
  the zeroth tensor power, \(\mathcal{L}^{\otimes 0} = \mathcal{O}_X\), and its
  global sections \(\Gamma(X, \mathcal{O}_X)\) form the degree-zero component of
  the section graded ring.
\end{lemma}



\begin{definition}\leanok
[Unit module of a scheme]
  \label{def:unitModule}
  \lean{AlgebraicGeometry.Scheme.Modules.unitModule}
  \uses{lem:moduleUnit_mathlib}
  For a scheme \(X\), the \emph{unit module} \(\mathbf{1}_X \in X.\Modules\) is the
  structure sheaf viewed as a module over itself, i.e.\ Mathlib's
  \(\mathrm{SheafOfModules.unit}\) (\cref{lem:moduleUnit_mathlib}) for
  \(\mathcal{O}_X\). It is the value of the zeroth tensor power,
  \(\mathcal{L}^{\otimes 0} = \mathbf{1}_X\) (\cref{def:sheafTensorPow}).
\end{definition}

\begin{definition}\leanok
[Tensor product of sheaves of modules]
  \label{def:sheafTensorObj}
  \lean{AlgebraicGeometry.Scheme.Modules.tensorObj}
  \uses{lem:presheafModule_monoidal_mathlib,
    lem:presheafModule_sheafification_mathlib, def:schemeModuleSheafification}
  % SOURCE: [Stacks Project], "Sheaves of Modules", \S "Tensor product" (Tag
  % 01CA), prose definition (read from references/stacks-modules.tex,
  % L2282-L2297).
  % SOURCE QUOTE:
  % "Let $(X, \mathcal{O}_X)$ be a ringed space and let
  %  $\mathcal{F}$ and $\mathcal{G}$ be $\mathcal{O}_X$-modules.
  %  We define first the {\it tensor product presheaf}
  %  $$
  %  \mathcal{F} \otimes_{p, \mathcal{O}_X} \mathcal{G}
  %  $$
  %  as the rule which assigns to $U \subset X$ open the
  %  $\mathcal{O}_X(U)$-module
  %  $\mathcal{F}(U) \otimes_{\mathcal{O}_X(U)} \mathcal{G}(U)$.
  %  Having defined this we define the {\it tensor product sheaf}
  %  as the sheafification of the above:
  %  $$
  %  \mathcal{F} \otimes_{\mathcal{O}_X} \mathcal{G}
  %  =
  %  (\mathcal{F} \otimes_{p, \mathcal{O}_X} \mathcal{G})^\#
  %  $$"
  \textit{Source: [Stacks Project], "Sheaves of Modules", Tag 01CA (tensor
  product).}
  Let \(\mathcal{F}, \mathcal{G}\) be sheaves of \(\mathcal{O}_X\)-modules. Their
  \emph{tensor product} is the sheaf of \(\mathcal{O}_X\)-modules
  \[
    \mathcal{F} \otimes_{\mathcal{O}_X} \mathcal{G}
      \;=\; \bigl(\mathcal{F} \otimes_{p,\mathcal{O}_X} \mathcal{G}\bigr)^{\#},
  \]
  the sheafification (\cref{lem:presheafModule_sheafification_mathlib}) of the
  tensor product presheaf \(\mathcal{F} \otimes_{p,\mathcal{O}_X} \mathcal{G}\),
  whose value on an open \(U\) is the \(\mathcal{O}_X(U)\)-module
  \(\mathcal{F}(U) \otimes_{\mathcal{O}_X(U)} \mathcal{G}(U)\). Concretely, the
  tensor product presheaf is the Mathlib presheaf-level monoidal product
  (\cref{lem:presheafModule_monoidal_mathlib}) of the underlying presheaves of
  modules of \(\mathcal{F}\) and \(\mathcal{G}\), and
  \(\mathcal{F} \otimes_{\mathcal{O}_X} \mathcal{G}\) is its sheafification. This
  construction is symmetric and associative because the presheaf tensor product
  is, and sheafification preserves the coherence isomorphisms.
\end{definition}

\begin{definition}\leanok
[Tensor power of a line bundle]
  \label{def:sheafTensorPow}
  \lean{AlgebraicGeometry.Scheme.Modules.tensorPow}
  \uses{def:sheafTensorObj, def:unitModule, lem:moduleUnit_mathlib}
  % SOURCE: [Stacks Project], "Sheaves of Modules", \S "Invertible modules",
  % `definition-powers' (Tag 01CU) (read from references/stacks-modules.tex,
  % L4220-L4228).
  % SOURCE QUOTE:
  % "Let $(X, \mathcal{O}_X)$ be a ringed space. Given an invertible sheaf
  %  $\mathcal{L}$ on $X$ and $n \in \mathbf{Z}$ we define the
  %  $n$th {\it tensor power} $\mathcal{L}^{\otimes n}$ of $\mathcal{L}$
  %  as the image of $\mathcal{O}_X$ under applying the equivalence
  %  $\mathcal{F} \mapsto\mathcal{F} \otimes_{\mathcal{O}_X} \mathcal{L}$
  %  exactly $n$ times."
  \textit{Source: [Stacks Project], "Sheaves of Modules", Tag 01CU (tensor
  powers).}
  Let \(\mathcal{L}\) be a sheaf of \(\mathcal{O}_X\)-modules (in the application
  an invertible sheaf / line bundle \(L_s\)). For \(m \in \mathbb{N}\) the
  \emph{\(m\)th tensor power} \(\mathcal{L}^{\otimes m}\) is defined by recursion
  on \(m\):
  \[
    \mathcal{L}^{\otimes 0} = \mathcal{O}_X
      \quad(\text{the unit module, \cref{lem:moduleUnit_mathlib}}),
    \qquad
    \mathcal{L}^{\otimes (m+1)}
      = \mathcal{L}^{\otimes m} \otimes_{\mathcal{O}_X} \mathcal{L}
        \quad(\cref{def:sheafTensorObj}).
  \]
  We restrict to non-negative powers, which is all the section graded ring uses;
  the Stacks definition extends \(\mathcal{L}^{\otimes n}\) to all
  \(n \in \mathbf{Z}\) using invertibility, but negative powers are not needed
  here. For \(\mathcal{L}\) invertible each \(\mathcal{L}^{\otimes m}\) is again
  invertible.
\end{definition}



\begin{definition}\leanok
[Twist of a coherent sheaf by a tensor power]
  \label{def:sheafModuleTwist}
  \lean{AlgebraicGeometry.Scheme.Modules.moduleTensorPow}
  \uses{def:sheafTensorObj, def:sheafTensorPow}
  Let \(\mathcal{F}\) be a sheaf of \(\mathcal{O}_X\)-modules (in the application
  a coherent sheaf \(\mathcal{F}_s\)) and \(\mathcal{L}\) a line bundle. For
  \(m \in \mathbb{N}\) the \emph{\(m\)-twist} of \(\mathcal{F}\) by
  \(\mathcal{L}\) is
  \[
    \mathcal{F}(m) \;:=\;
      \mathcal{F} \otimes_{\mathcal{O}_X} \mathcal{L}^{\otimes m}
      \qquad(\cref{def:sheafTensorObj}, \cref{def:sheafTensorPow}),
  \]
  the degree-\(m\) carrier of the section graded module. By construction
  \(\mathcal{F}(0) = \mathcal{F} \otimes_{\mathcal{O}_X} \mathcal{O}_X \cong
  \mathcal{F}\).
\end{definition}


The structural isomorphisms below --- the left unitor and the braiding --- are
the parts of the (would-be) symmetric monoidal structure on \(X.\Modules\) that
descend through sheafification from \cref{lem:presheafModule_monoidal_mathlib}
using only \emph{functoriality} of the sheafification functor
(\cref{def:schemeModuleSheafification}) and, for the unitor, the reflective
counit isomorphism --- no strong monoidality of sheafification is required. They
are, together with the associator \cref{cor:sheafTensorObjAssoc}, the structural
inputs of the tensor-power comparison \cref{lem:sheafTensorPow_add}.

\begin{definition}\leanok
[Left unitor of the sheaf tensor product]
  \label{def:tensorObjUnitIso}
  \lean{AlgebraicGeometry.Scheme.Modules.tensorObjUnitIso}
  \uses{def:sheafTensorObj, def:unitModule, def:schemeModuleSheafification,
    lem:presheafModule_monoidal_mathlib,
    lem:presheafModule_sheafification_mathlib}
  For a sheaf of \(\mathcal{O}_X\)-modules \(\mathcal{G}\) the \emph{left unitor}
  is the isomorphism
  \[
    \lambda_{\mathcal{G}} :
      \mathbf{1}_X \otimes_{\mathcal{O}_X} \mathcal{G}
      \;\xrightarrow{\ \sim\ }\; \mathcal{G}
  \]
  of sheaves of modules, where \(\mathbf{1}_X\) is the unit module
  (\cref{def:unitModule}). It is the sheafification
  (\cref{def:schemeModuleSheafification}) of the presheaf-level left unitor
  \(\lambda^p\) of \cref{lem:presheafModule_monoidal_mathlib}, composed with the
  counit isomorphism of the sheafification adjunction that identifies the
  sheafification of the underlying presheaf of a sheaf of modules with the sheaf
  itself. It is the base case (\(m = 0\)) of the tensor-power comparison
  \cref{lem:sheafTensorPow_add}.
\end{definition}

\begin{definition}\leanok
[Braiding of the sheaf tensor product]
  \label{def:tensorBraiding}
  \lean{AlgebraicGeometry.Scheme.Modules.tensorBraiding}
  \uses{def:sheafTensorObj, def:schemeModuleSheafification,
    lem:presheafModule_monoidal_mathlib}
  For sheaves of \(\mathcal{O}_X\)-modules \(\mathcal{F}, \mathcal{G}\) the
  \emph{braiding} is the symmetry isomorphism
  \[
    \beta_{\mathcal{F},\mathcal{G}} :
      \mathcal{F} \otimes_{\mathcal{O}_X} \mathcal{G}
      \;\xrightarrow{\ \sim\ }\;
      \mathcal{G} \otimes_{\mathcal{O}_X} \mathcal{F}
  \]
  of sheaves of modules, obtained by applying the sheafification functor
  (\cref{def:schemeModuleSheafification}) to the symmetric braiding of the
  presheaf monoidal structure (\cref{lem:presheafModule_monoidal_mathlib}). Being
  pure sheafification-functoriality of an isomorphism, it requires no monoidal
  structure on \(X.\Modules\); it is the symmetry used in the inductive step of
  \cref{lem:sheafTensorPow_add}.
\end{definition}

\begin{definition}

\leanok
[Right unitor of the sheaf tensor product]
  \label{def:tensorObjRightUnitor}
  \lean{AlgebraicGeometry.Scheme.Modules.tensorObjRightUnitor}
  \uses{def:sheafTensorObj, def:unitModule, def:tensorBraiding,
    def:tensorObjUnitIso}
  For a sheaf of \(\mathcal{O}_X\)-modules \(\mathcal{G}\) the \emph{right unitor}
  is the isomorphism
  \[
    \rho_{\mathcal{G}} :
      \mathcal{G} \otimes_{\mathcal{O}_X} \mathbf{1}_X
      \;\xrightarrow{\ \sim\ }\; \mathcal{G}
  \]
  of sheaves of modules, where \(\mathbf{1}_X\) is the unit module
  (\cref{def:unitModule}). It is defined as the composite of the braiding that
  swaps the two factors with the left unitor,
  \[
    \mathcal{G} \otimes_{\mathcal{O}_X} \mathbf{1}_X
      \;\xrightarrow{\ \beta_{\mathcal{G}, \mathbf{1}_X}\ }\;
      \mathbf{1}_X \otimes_{\mathcal{O}_X} \mathcal{G}
      \;\xrightarrow{\ \lambda_{\mathcal{G}}\ }\;
      \mathcal{G},
    \qquad
    \rho_{\mathcal{G}}
      = \lambda_{\mathcal{G}} \circ \beta_{\mathcal{G}, \mathbf{1}_X},
  \]
  using the braiding \(\beta\) (\cref{def:tensorBraiding}) and the left unitor
  \(\lambda\) (\cref{def:tensorObjUnitIso}). Like those two, it descends from the
  presheaf symmetric monoidal structure through sheafification-functoriality and
  the reflective counit, requiring no monoidal structure on \(X.\Modules\). It is
  the base of the right-unitality reduction of the tensor-power comparison
  (\cref{lem:tensorPowAdd_rightUnit}).
\end{definition}





The strong-monoidality comparison \cref{lem:isIso_sheafification_whiskerRight_unit}
below is assembled from two reduction lemmas --- a localization criterion
(\cref{lem:isIso_sheafification_map_iff}) and the local-isomorphism property of
the unit (\cref{lem:localIso_toPresheaf_map_unit}) --- together with the
coequalizer presentation of the relative module-presheaf tensor product, the single
Mathlib-absent brick on which the comparison rests.

\begin{lemma}\leanok
[Localization criterion for module sheafification]
  \label{lem:isIso_sheafification_map_iff}
  \lean{AlgebraicGeometry.Scheme.Modules.isIso_sheafification_map_iff}
  \uses{def:schemeModuleSheafification, lem:presheafModule_sheafification_mathlib}
  Let \(f : \mathcal{P} \to \mathcal{Q}\) be a morphism of presheaves of
  \(\mathcal{O}_X\)-modules. Then the sheafified morphism \(f^{\#}\)
  (\cref{def:schemeModuleSheafification}) is an isomorphism of sheaves of modules
  if and only if the underlying morphism of abelian-group presheaves of \(f\)
  lies in the class \(J.W\) of local isomorphisms (weak equivalences) for the
  opens topology \(J\) on \(X\), i.e.\ \(f\) lies in the inverse image of \(J.W\)
  along the forgetful functor to abelian-group presheaves.
\end{lemma}

\begin{proof}
  \uses{lem:presheafModule_sheafification_mathlib}
  \leanok
  Module sheafification (\cref{def:schemeModuleSheafification}) is the localization
  functor at the class \(W\) of morphisms of presheaves of modules whose underlying
  abelian-presheaf morphism lies in \(J.W\), i.e.\ the inverse image of \(J.W\) along
  the forgetful functor. The localization criterion for a reflective adjunction
  (\cref{lem:presheafModule_sheafification_mathlib}) then gives the stated
  equivalence.
\end{proof}


\begin{lemma}\leanok
[The sheafification unit is an abelian local isomorphism]
  \label{lem:localIso_toPresheaf_map_unit}
  \lean{AlgebraicGeometry.Scheme.Modules.localIso_toPresheaf_map_unit}
  \uses{def:schemeModuleSheafification, lem:presheafModule_sheafification_mathlib}
  For a presheaf of \(\mathcal{O}_X\)-modules \(\mathcal{P}\), the underlying
  morphism of abelian-group presheaves of the sheafification unit
  \(\eta_{\mathcal{P}} : \mathcal{P} \to \mathcal{P}^{\#}\)
  (\cref{def:schemeModuleSheafification}) is, by construction, the unit of the
  abelian associated-sheaf adjunction for the opens topology \(J\) on \(X\); it is
  therefore a local isomorphism and lies in the class \(J.W\).
\end{lemma}

\begin{proof}
  \uses{lem:presheafModule_sheafification_mathlib}
  \leanok
  The forgetful functor to abelian-group presheaves carries the module
  sheafification unit to the unit of the abelian associated-sheaf adjunction, which
  is a local isomorphism, hence a member of \(J.W\).
\end{proof}













The coequalizer presentation of the relative module-presheaf tensor product is
the single Mathlib-absent brick on which the strong-monoidality comparison rests.
It is built from the abelian-group universal property of the relative tensor
product (\cref{lem:tensorProduct_liftAddHom_mathlib}), promoted to the presheaf
category where colimits are computed objectwise
(\cref{lem:evaluationJointlyReflectsColimits_mathlib}) and the apex is identified
with the presheaf tensor product (\cref{lem:presheaf_tensorObj_obj_mathlib}).

\begin{lemma}[Universal property of the relative tensor product]
  \label{lem:tensorProduct_liftAddHom_mathlib}
  \lean{TensorProduct.liftAddHom}
  \mathlibok
  \textit{Provided by Mathlib (\texttt{Mathlib.LinearAlgebra.TensorProduct.Basic}).}
  Let \(S\) be a commutative ring and \(M, N, P\) additive commutative groups with
  \(M, N\) modules over \(S\). An additive map \(f : M \to (N \to_{+} P)\) that is
  \(S\)-balanced (\(f(s \cdot m)(n) = f(m)(s \cdot n)\)) factors uniquely through
  the canonical surjection \(M \otimes_{\mathbb{Z}} N \to M \otimes_S N\),
  producing a homomorphism \(M \otimes_S N \to P\). This is the abelian-group
  universal property of the relative tensor product.
\end{lemma}

\begin{lemma}[Evaluation jointly reflects colimits]
  \label{lem:evaluationJointlyReflectsColimits_mathlib}
  \lean{CategoryTheory.Limits.evaluationJointlyReflectsColimits}
  \mathlibok
  \textit{Provided by Mathlib
  (\texttt{Mathlib.CategoryTheory.Limits.FunctorCategory.Basic}).}
  Let \(F : \mathcal{J} \to (\mathcal{K} \to \mathcal{C})\) be a diagram in a
  functor category. A cocone over \(F\) is a colimit as soon as its image under
  every evaluation functor \(\mathrm{ev}_U : (\mathcal{K} \to \mathcal{C}) \to
  \mathcal{C}\), \(U \in \mathcal{K}\), is a colimit. In particular colimits in a
  functor category are computed objectwise.
\end{lemma}

\begin{lemma}[Objectwise value of the presheaf tensor product]
  \label{lem:presheaf_tensorObj_obj_mathlib}
  \lean{PresheafOfModules.Monoidal.tensorObj_obj}
  \mathlibok
  \textit{Provided by Mathlib
  (\texttt{Mathlib.Algebra.Category.ModuleCat.Presheaf.Monoidal}).}
  For presheaves of modules \(\mathcal{M}_1, \mathcal{M}_2\) over a presheaf of
  rings and an object \(U\), the value of the presheaf monoidal tensor product is
  \[
    (\mathcal{M}_1 \otimes_p \mathcal{M}_2)(U)
      = \mathcal{M}_1(U) \otimes_{\mathcal{O}_X(U)} \mathcal{M}_2(U),
  \]
  the objectwise tensor product of the values over the ring \(\mathcal{O}_X(U)\).
\end{lemma}

\begin{definition}\leanok
[Objectwise \(\mathbb{Z}\)-tensor presheaf]
  \label{def:relTensorDomainPresheaf}
  \lean{AlgebraicGeometry.Scheme.Modules.relTensorDomainPresheaf}
  For presheaves of \(\mathcal{O}_X\)-modules \(\mathcal{P}, \mathcal{Q}\), the
  \emph{objectwise \(\mathbb{Z}\)-tensor presheaf} is the functor
  \[
    U \;\longmapsto\; \mathcal{P}(U) \otimes_{\mathbb{Z}} \mathcal{Q}(U)
      \;:\; (\mathrm{Opens}\,X)^{\mathrm{op}} \to \mathrm{Ab},
  \]
  with restriction maps the \(\mathbb{Z}\)-tensors of the underlying restriction
  maps of \(\mathcal{P}\) and \(\mathcal{Q}\). It is the apex-adjacent presheaf of
  the relative-tensor cofork; no objectwise \(\mathbb{Z}\)-tensor of
  abelian-group presheaves is supplied by Mathlib.
\end{definition}

\begin{definition}\leanok
[Triple \(\mathbb{Z}\)-tensor presheaf]
  \label{def:relTensorTriplePresheaf}
  \lean{AlgebraicGeometry.Scheme.Modules.relTensorTriplePresheaf}
  For presheaves of \(\mathcal{O}_X\)-modules \(\mathcal{P}, \mathcal{Q}\), the
  \emph{triple \(\mathbb{Z}\)-tensor presheaf} is the functor
  \[
    U \;\longmapsto\; \mathcal{P}(U) \otimes_{\mathbb{Z}}
      \bigl(\mathcal{O}_X(U) \otimes_{\mathbb{Z}} \mathcal{Q}(U)\bigr)
      \;:\; (\mathrm{Opens}\,X)^{\mathrm{op}} \to \mathrm{Ab},
  \]
  the middle factor restricting along the ring restriction maps of
  \(\mathcal{O}_X\). It is the domain row of the relative-tensor coequalizer
  presentation \cref{lem:relativeTensor_as_coequalizer}.
\end{definition}

\begin{definition}\leanok
[Left-action row of the coequalizer]
  \label{def:relTensorActL}
  \lean{AlgebraicGeometry.Scheme.Modules.relTensorActL}
  \uses{def:relTensorTriplePresheaf, def:relTensorDomainPresheaf}
  The \emph{left-action} natural transformation
  \(a_L : \cref{def:relTensorTriplePresheaf} \to \cref{def:relTensorDomainPresheaf}\)
  has component at \(U\) the map collapsing the middle ring factor through the
  scalar action of \(\mathcal{O}_X(U)\) on \(\mathcal{P}(U)\),
  \(m \otimes (s \otimes n) \mapsto (s \cdot m) \otimes n\). Naturality in \(U\)
  is the compatibility of the module action with the restriction maps.
\end{definition}

\begin{definition}\leanok
[Right-action row of the coequalizer]
  \label{def:relTensorActR}
  \lean{AlgebraicGeometry.Scheme.Modules.relTensorActR}
  \uses{def:relTensorTriplePresheaf, def:relTensorDomainPresheaf}
  The \emph{right-action} natural transformation
  \(a_R : \cref{def:relTensorTriplePresheaf} \to \cref{def:relTensorDomainPresheaf}\)
  has component at \(U\) the map collapsing the middle ring factor through the
  scalar action of \(\mathcal{O}_X(U)\) on \(\mathcal{Q}(U)\),
  \(m \otimes (s \otimes n) \mapsto m \otimes (s \cdot n)\). Naturality is the
  compatibility of the module action with the restriction maps.
\end{definition}

\begin{definition}\leanok
[Projection cofork leg]
  \label{def:relTensorProj}
  \lean{AlgebraicGeometry.Scheme.Modules.relTensorProj}
  \uses{def:relTensorDomainPresheaf, lem:presheaf_tensorObj_obj_mathlib}
  The \emph{projection} natural transformation
  \(\pi : \cref{def:relTensorDomainPresheaf} \to (\mathcal{P} \otimes_p
  \mathcal{Q})\) has component at \(U\) the canonical quotient
  \(\mathcal{P}(U) \otimes_{\mathbb{Z}} \mathcal{Q}(U) \to \mathcal{P}(U)
  \otimes_{\mathcal{O}_X(U)} \mathcal{Q}(U)\) onto the relative tensor, the apex of
  the cofork, identified with the value of the presheaf monoidal tensor product
  (\cref{lem:presheaf_tensorObj_obj_mathlib}).
\end{definition}

\begin{lemma}\leanok
[Cofork condition for the relative-tensor coequalizer]
  \label{lem:snap_relTensorActL_proj_eq}
  \lean{AlgebraicGeometry.Scheme.Modules.relTensorActL_proj_eq}
  \uses{def:relTensorActL, def:relTensorActR, def:relTensorProj,
    lem:relativeTensor_objectwise_coequalizer}
  The left- and right-action rows compose equally with the projection,
  \[
    a_L \circ \pi = a_R \circ \pi,
  \]
  as natural transformations of \((\mathrm{Opens}\,X)^{\mathrm{op}} \to
  \mathrm{Ab}\) (\cref{def:relTensorActL}, \cref{def:relTensorActR},
  \cref{def:relTensorProj}).
\end{lemma}

\begin{proof}
  \uses{lem:relativeTensor_objectwise_coequalizer}
  \leanok
  Objectwise at each open \(U\) the identity is the cofork condition of the
  objectwise coequalizer \cref{lem:relativeTensor_objectwise_coequalizer}: both
  composites send \(m \otimes (s \otimes n)\) to the relative-tensor class
  \((s \cdot m) \otimes_{\mathcal{O}_X(U)} n = m \otimes_{\mathcal{O}_X(U)}
  (s \cdot n)\), equal by the defining balancing relation. Naturality makes this
  an equality of natural transformations.
\end{proof}

\begin{lemma}\leanok
[\(\mathbb{Z}\)-whiskered local isomorphisms are local isomorphisms]
  \label{lem:snap_ztensor_whisker_localIso}
  \lean{AlgebraicGeometry.Scheme.Modules.ztensor_whisker_localIso}
  \uses{lem:relativeTensor_as_coequalizer,
    lem:relativeTensor_objectwise_coequalizer, def:relTensorTriplePresheaf,
    def:relTensorDomainPresheaf}
  Let \(f : \mathcal{P} \to \mathcal{Q}\) be a morphism of presheaves of
  \(\mathcal{O}_X\)-modules whose underlying abelian-presheaf morphism lies in the
  class \(J.W\) of local isomorphisms for the opens topology \(J\) on \(X\). Then
  for any presheaf of modules \(\mathcal{R}\) the underlying abelian morphism of
  the right-whiskered map \(f \rhd \mathcal{R} : \mathcal{P} \otimes_p \mathcal{R}
  \to \mathcal{Q} \otimes_p \mathcal{R}\) (in the presheaf monoidal structure
  \cref{lem:presheafModule_monoidal_mathlib}) again lies in \(J.W\).
\end{lemma}

\begin{proof}
  \uses{def:relTensorDomainPresheaf, def:relTensorTriplePresheaf,
    lem:relativeTensor_as_coequalizer, lem:relativeTensor_objectwise_coequalizer}
    \leanok
  Two reductions.

  \emph{(i) \(\mathbb{Z}\)-tensor whiskering preserves \(J.W\).} The category
  of abelian-group presheaves on \(X\) embeds, by a carrier-preserving
  equivalence, into presheaves valued in modules over a (universe lift of)
  \(\mathbb{Z}\). The latter value category is braided monoidal closed, so the
  pointwise (Day) monoidal structure on its presheaf category satisfies: right
  whiskering by a fixed object carries \(J.W\) into \(J.W\) (the reflection
  argument for monoidal sites — sheafification is a monoidal localization for a
  braided closed value category). An equivalence of value categories is a left
  adjoint in both directions, hence preserves sheafification both ways and
  identifies the two classes \(J.W\); transporting along it shows that the
  \(\mathbb{Z}\)-tensor whiskering \(g \otimes_{\mathbb{Z}} \mathrm{id}\) of any
  \(g \in J.W\) by a fixed abelian presheaf again lies in \(J.W\).

  \emph{(ii) Descent through the coequalizer presentation.} By
  \cref{lem:relativeTensor_as_coequalizer}, the underlying abelian presheaf of
  \(\mathcal{P} \otimes_p \mathcal{R}\) is the coequalizer (computed objectwise)
  of the two action rows out of the triple \(\mathbb{Z}\)-tensor presheaf
  (\cref{def:relTensorTriplePresheaf}), and the underlying abelian morphism of
  \(f \rhd \mathcal{R}\) is the map induced on coequalizer points by
  \(\mathbb{Z}\)-whiskering the two rows with \(f\). Abelian sheafification is
  a left adjoint, hence preserves these coequalizers; by (i) it inverts the
  whiskered rows, hence it inverts the induced map on coequalizer points. A
  morphism inverted by sheafification lies in \(J.W\), which is the claim.
\end{proof}



\begin{lemma}\leanok
[Relative tensor product as an objectwise coequalizer]
  \label{lem:relativeTensor_objectwise_coequalizer}
  \lean{AlgebraicGeometry.Scheme.Modules.RelativeTensorCoequalizer.isColimitCofork}
  \uses{lem:tensorProduct_liftAddHom_mathlib}
  Let \(S\) be a commutative ring and \(M, N\) two \(S\)-modules. Then the relative
  tensor product \(M \otimes_S N\) is the coequalizer, in the category of abelian
  groups, of the two \(S\)-action maps
  \[
    M \otimes_{\mathbb{Z}} (S \otimes_{\mathbb{Z}} N)
      \;\overset{a_L,\,a_R}{\rightrightarrows}\;
      M \otimes_{\mathbb{Z}} N,
    \qquad
    a_L : m \otimes s \otimes n \mapsto (s \cdot m) \otimes n,
    \quad
    a_R : m \otimes s \otimes n \mapsto m \otimes (s \cdot n),
  \]
  with cofork leg the canonical projection \(\pi : M \otimes_{\mathbb{Z}} N \to
  M \otimes_S N\). This is the objectwise content of
  \cref{lem:relativeTensor_as_coequalizer}.
\end{lemma}

\begin{proof}
  \uses{lem:tensorProduct_liftAddHom_mathlib}
  \leanok
  Write \(a_N : S \otimes_{\mathbb{Z}} N \to N\) and
  \(a_M : M \otimes_{\mathbb{Z}} S \to M\) for the \(\mathbb{Z}\)-linear maps
  induced by the scalar action (\(s \otimes n \mapsto s \cdot n\),
  \(m \otimes s \mapsto s \cdot m\)); the right action map is
  \(a_R = \mathrm{id}_M \otimes a_N\) and the left action map is
  \(a_L = (a_M \otimes \mathrm{id}_N) \circ \mathrm{assoc}^{-1}\), where
  \(\mathrm{assoc}\) is the \(\mathbb{Z}\)-tensor associator. The projection
  \(\pi\) is the additive homomorphism produced by the abelian universal property
  \cref{lem:tensorProduct_liftAddHom_mathlib} from the \(S\)-balanced bilinear map
  \((m, n) \mapsto m \otimes_S n\).

  \emph{Cofork condition.} On an elementary tensor
  \(m \otimes (s \otimes n)\) we have
  \(\pi(a_L(m \otimes s \otimes n)) = (s \cdot m) \otimes_S n
  = m \otimes_S (s \cdot n) = \pi(a_R(m \otimes s \otimes n))\), the middle
  equality being the defining \(S\)-balancing relation of \(M \otimes_S N\); since
  elementary tensors generate, \(\pi \circ a_L = \pi \circ a_R\).

  \emph{Universal property.} Given any abelian group \(P\) and a homomorphism
  \(g : M \otimes_{\mathbb{Z}} N \to P\) with \(g \circ a_L = g \circ a_R\), the
  bilinear map \((m, n) \mapsto g(m \otimes n)\) is \(S\)-balanced (that is exactly
  the cofork condition read on elementary tensors), so by
  \cref{lem:tensorProduct_liftAddHom_mathlib} it factors as a unique
  \(\widetilde{g} : M \otimes_S N \to P\) with \(\widetilde{g} \circ \pi = g\).
  Existence is this descended map; uniqueness follows because \(\pi\) is
  surjective, hence an epimorphism of abelian groups, so \(\widetilde{g} \circ \pi
  = g\) determines \(\widetilde{g}\). Packaging existence and uniqueness gives the
  colimit cofork.

\end{proof}


\begin{lemma}\leanok
[Relative tensor product as a presheaf coequalizer]
  \label{lem:relativeTensor_as_coequalizer}
  \lean{AlgebraicGeometry.Scheme.Modules.relativeTensorCoequalizerIso}
  \uses{lem:relativeTensor_objectwise_coequalizer, def:relTensorActL,
    def:relTensorActR, def:relTensorProj, lem:snap_relTensorActL_proj_eq,
    lem:evaluationJointlyReflectsColimits_mathlib,
    lem:presheaf_tensorObj_obj_mathlib, lem:presheafModule_monoidal_mathlib}
  For presheaves of \(\mathcal{O}_X\)-modules \(\mathcal{P}, \mathcal{Q}\), the
  underlying abelian-group presheaf of the presheaf-level relative tensor product
  \(\mathcal{P} \otimes_p \mathcal{Q}\) (\cref{lem:presheafModule_monoidal_mathlib})
  is the coequalizer, in the functor category
  \((\mathrm{Opens}\,X)^{\mathrm{op}} \to \mathrm{Ab}\), of the parallel pair of
  action rows \(a_L\) (\cref{def:relTensorActL}), \(a_R\)
  (\cref{def:relTensorActR}) with cofork leg the projection \(\pi\)
  (\cref{def:relTensorProj}). It is the presheaf-level promotion of the objectwise
  coequalizer \cref{lem:relativeTensor_objectwise_coequalizer}.
\end{lemma}

\begin{proof}
  \uses{lem:presheafModule_monoidal_mathlib,
    lem:relativeTensor_objectwise_coequalizer,
    lem:evaluationJointlyReflectsColimits_mathlib,
    lem:presheaf_tensorObj_obj_mathlib,
    def:relTensorActL, def:relTensorActR, def:relTensorProj,
    lem:snap_relTensorActL_proj_eq}
    \leanok
  The argument promotes the objectwise coequalizer of
  \cref{lem:relativeTensor_objectwise_coequalizer} to the functor category, in
  three steps.

  \textbf{1. Objectwise.} For each open \(U\), apply
  \cref{lem:relativeTensor_objectwise_coequalizer} with \(S = \mathcal{O}_X(U)\),
  \(M = \mathcal{P}(U)\), \(N = \mathcal{Q}(U)\): the abelian group
  \(\mathcal{P}(U) \otimes_{\mathcal{O}_X(U)} \mathcal{Q}(U)\) is the coequalizer
  in \(\mathrm{AddCommGrp}\) of the two \(\mathcal{O}_X(U)\)-action maps
  \[
    \mathcal{P}(U) \otimes_{\mathbb{Z}} \mathcal{O}_X(U)
      \otimes_{\mathbb{Z}} \mathcal{Q}(U)
    \;\rightrightarrows\;
    \mathcal{P}(U) \otimes_{\mathbb{Z}} \mathcal{Q}(U),
  \]
  with cofork map the projection \(\pi_U\).

  \textbf{2. Naturality / promotion.} The three objectwise families are precisely
  the natural transformations of presheaves on \(X\)
  \[
    \mathcal{P} \otimes_{\mathbb{Z}} \mathcal{O}_X
      \otimes_{\mathbb{Z}} \mathcal{Q}
    \;\overset{a_L,\,a_R}{\rightrightarrows}\;
    \mathcal{P} \otimes_{\mathbb{Z}} \mathcal{Q}
    \;\overset{\pi}{\longrightarrow}\;
    \bigl(U \mapsto \mathcal{P}(U) \otimes_{\mathcal{O}_X(U)} \mathcal{Q}(U)\bigr)
  \]
  built in the previous definitions: the parallel pair is the left- and
  right-action transformations \(a_L\) (\cref{def:relTensorActL}) and \(a_R\)
  (\cref{def:relTensorActR}) from the triple presheaf
  (\cref{def:relTensorTriplePresheaf}) to the \(\mathbb{Z}\)-tensor presheaf
  (\cref{def:relTensorDomainPresheaf}), and the cofork leg is the projection
  \(\pi\) (\cref{def:relTensorProj}) onto the apex. These are morphisms of
  functors \((\mathrm{Opens}\,X)^{\mathrm{op}} \to \mathrm{Ab}\); naturality of
  each in \(U\) is the compatibility of the module action with the restriction
  maps of \(\mathcal{P}, \mathcal{O}_X, \mathcal{Q}\) (the restriction of
  \(s \cdot m\) is \((s|_V) \cdot (m|_V)\), and likewise for the projection),
  recorded in those definitions. Objectwise at each \(U\) the cofork
  \((a_L, a_R, \pi)_U\) is the colimit cofork of step 1. Colimits in a functor
  category are computed objectwise --- by
  \cref{lem:evaluationJointlyReflectsColimits_mathlib} a cocone is a colimit as
  soon as every evaluation \((\mathrm{ev}_U)\) of it is --- so the objectwise
  coequalizers of step 1 promote to a coequalizer of the displayed parallel pair
  \(a_L, a_R\) with cofork \(\pi\) in
  \((\mathrm{Opens}\,X)^{\mathrm{op}} \to \mathrm{Ab}\).

  \textbf{3. Identification of the apex.} The apex presheaf
  \(U \mapsto \mathcal{P}(U) \otimes_{\mathcal{O}_X(U)} \mathcal{Q}(U)\) is, by the
  objectwise tensor formula of the presheaf monoidal structure
  (\cref{lem:presheaf_tensorObj_obj_mathlib}: \((\mathcal{M}_1 \otimes_p
  \mathcal{M}_2)(U) = \mathcal{M}_1(U) \otimes_{\mathcal{O}_X(U)}
  \mathcal{M}_2(U)\)), exactly the underlying abelian-group presheaf of the
  relative tensor product \(\mathcal{P} \otimes_{p,\mathcal{O}_X} \mathcal{Q}\)
  (\cref{lem:presheafModule_monoidal_mathlib}). Transporting the colimit of step 2
  along this identification exhibits
  \(\mathcal{P} \otimes_{p,\mathcal{O}_X} \mathcal{Q}\), as an abelian-group
  presheaf, as the coequalizer of the parallel pair, naturally in \(U\).
\end{proof}




\begin{lemma}\leanok
[Sheafification inverts the whiskered localization unit]
  \label{lem:isIso_sheafification_whiskerRight_unit}
  \lean{AlgebraicGeometry.Scheme.Modules.isIso_sheafification_whiskerRight_unit}
  \uses{def:schemeModuleSheafification, lem:isIso_sheafification_map_iff,
    lem:localIso_toPresheaf_map_unit, lem:snap_ztensor_whisker_localIso,
    lem:presheafModule_monoidal_mathlib,
    lem:presheafModule_sheafification_mathlib}
  For presheaves of \(\mathcal{O}_X\)-modules \(\mathcal{P}, \mathcal{Q}\), let
  \(\eta_{\mathcal{P}} : \mathcal{P} \to \mathcal{P}^{\#}\) be the unit of the
  sheafification adjunction (\cref{def:schemeModuleSheafification}). The
  sheafification of the right-whiskered map
  \(\eta_{\mathcal{P}} \rhd \mathcal{Q} : \mathcal{P} \otimes_p \mathcal{Q} \to
  \mathcal{P}^{\#} \otimes_p \mathcal{Q}\) (in the presheaf monoidal structure
  \cref{lem:presheafModule_monoidal_mathlib}), namely
  \[
    (\eta_{\mathcal{P}} \rhd \mathcal{Q})^{\#} :
      (\mathcal{P} \otimes_p \mathcal{Q})^{\#}
      \;\xrightarrow{\ \sim\ }\;
      (\mathcal{P}^{\#} \otimes_p \mathcal{Q})^{\#},
  \]
  is an isomorphism of sheaves of modules. This is the strong-monoidality
  comparison of module sheafification on a whiskered unit --- the single
  Mathlib-absent brick for the sheaf-level associator
  (\cref{cor:sheafTensorObjAssoc}) and the tensor-power comparison
  \cref{lem:sheafTensorPow_add}.
\end{lemma}

\begin{proof}
  \uses{lem:presheafModule_monoidal_mathlib,
    lem:presheafModule_sheafification_mathlib,
    lem:isIso_sheafification_map_iff, lem:localIso_toPresheaf_map_unit,
    lem:relativeTensor_as_coequalizer,
    lem:relativeTensor_objectwise_coequalizer,
    lem:snap_ztensor_whisker_localIso, def:relTensorDomainPresheaf}
    \leanok
  Module sheafification is a localization functor: it inverts exactly the class
  \(W\) of morphisms of presheaves of modules whose underlying morphism of
  abelian-group presheaves is a local isomorphism --- that is, lies in the class
  \(J.W\) of weak equivalences attached to the Grothendieck topology \(J\) on
  \(X\) (the inverse image of \(J.W\) along the forgetful functor to
  abelian-group presheaves). By the localization criterion for a reflective
  adjunction, \((f)^{\#}\) is an isomorphism if and only if \(f \in W\), i.e.\
  if and only if the underlying abelian-presheaf morphism of \(f\) lies in
  \(J.W\). It therefore suffices to prove that the underlying abelian morphism
  of \(\eta_{\mathcal{P}} \rhd \mathcal{Q}\) lies in \(J.W\).

  \textbf{The unit is a local isomorphism.} The sheafification unit
  \(\eta_{\mathcal{P}} : \mathcal{P} \to (\mathcal{P})^{\#}\) has underlying
  abelian-presheaf morphism in \(J.W\)
  (\cref{lem:localIso_toPresheaf_map_unit} — the defining property of the
  associated-sheaf unit for the topology \(J\)).

  \textbf{Whiskering preserves \(J.W\).} By
  \cref{lem:snap_ztensor_whisker_localIso}, right-whiskering a morphism of
  presheaves of modules whose underlying abelian morphism lies in \(J.W\) by
  the fixed presheaf \(\mathcal{Q}\) yields a morphism whose underlying
  abelian morphism again lies in \(J.W\) (proved there by descending the
  pointwise monoidal-whiskering statement through the relative-tensor
  coequalizer presentation of
  \cref{lem:relativeTensor_as_coequalizer}). Applied to
  \(\eta_{\mathcal{P}}\), the underlying abelian morphism of
  \(\eta_{\mathcal{P}} \rhd \mathcal{Q}\) lies in \(J.W\), so
  \(\eta_{\mathcal{P}} \rhd \mathcal{Q} \in W\), and
  \((\eta_{\mathcal{P}} \rhd \mathcal{Q})^{\#}\) is an isomorphism.
\end{proof}


\begin{corollary}\leanok
[Associator of the sheaf tensor product]
  \label{cor:sheafTensorObjAssoc}
  \lean{AlgebraicGeometry.Scheme.Modules.tensorObjAssoc}
  \uses{def:sheafTensorObj, def:schemeModuleSheafification,
    lem:isIso_sheafification_whiskerRight_unit,
    lem:presheafModule_monoidal_mathlib,
    lem:presheafModule_sheafification_mathlib}
  For sheaves of \(\mathcal{O}_X\)-modules \(\mathcal{A}, \mathcal{B},
  \mathcal{C}\) there is a canonical \emph{associator} isomorphism
  \[
    \alpha_{\mathcal{A},\mathcal{B},\mathcal{C}} :
      (\mathcal{A} \otimes_{\mathcal{O}_X} \mathcal{B}) \otimes_{\mathcal{O}_X}
        \mathcal{C}
      \;\xrightarrow{\ \sim\ }\;
      \mathcal{A} \otimes_{\mathcal{O}_X} (\mathcal{B} \otimes_{\mathcal{O}_X}
        \mathcal{C})
  \]
  of sheaves of modules, natural in its three arguments. It descends from the
  presheaf associator (\cref{lem:presheafModule_monoidal_mathlib}) through
  sheafification (\cref{def:schemeModuleSheafification}), the inner
  sheafifications of the two iterated tensor products (\cref{def:sheafTensorObj})
  being cleared by the strong-monoidality comparison
  \cref{lem:isIso_sheafification_whiskerRight_unit}.
\end{corollary}

\begin{proof}
  \uses{def:sheafTensorObj, lem:isIso_sheafification_whiskerRight_unit,
    lem:presheafModule_monoidal_mathlib,
    lem:presheafModule_sheafification_mathlib}
    \leanok
  Write \(a, b, c\) for the underlying presheaves of modules of
  \(\mathcal{A}, \mathcal{B}, \mathcal{C}\) and \((-)^{\#}\) for sheafification
  (\cref{lem:presheafModule_sheafification_mathlib}). By
  \cref{def:sheafTensorObj} the two iterated products unfold to
  \[
    (\mathcal{A} \otimes \mathcal{B}) \otimes \mathcal{C}
      = \bigl((a \otimes_p b)^{\#} \otimes_p c\bigr)^{\#},
    \qquad
    \mathcal{A} \otimes (\mathcal{B} \otimes \mathcal{C})
      = \bigl(a \otimes_p (b \otimes_p c)^{\#}\bigr)^{\#},
  \]
  each carrying an \emph{inner} sheafification that obstructs a direct appeal to
  the presheaf associator. The whiskered units remove these inner
  sheafifications. Applying \cref{lem:isIso_sheafification_whiskerRight_unit}
  with \(\mathcal{P} = a \otimes_p b\) and \(\mathcal{Q} = c\), the morphism
  \[
    (\eta_{a \otimes_p b} \rhd c)^{\#} :
    \bigl((a \otimes_p b) \otimes_p c\bigr)^{\#}
    \xrightarrow{\ \sim\ }
    \bigl((a \otimes_p b)^{\#} \otimes_p c\bigr)^{\#}
    = (\mathcal{A} \otimes \mathcal{B}) \otimes \mathcal{C}
  \]
  is an isomorphism; and, reading the same lemma through the braiding of the
  presheaf monoidal structure (\cref{lem:presheafModule_monoidal_mathlib}) so as
  to whisker the unit \(\eta_{b \otimes_p c}\) on the left of \(a\), the
  morphism
  \[
    (a \lhd \eta_{b \otimes_p c})^{\#} :
    \bigl(a \otimes_p (b \otimes_p c)\bigr)^{\#}
    \xrightarrow{\ \sim\ }
    \bigl(a \otimes_p (b \otimes_p c)^{\#}\bigr)^{\#}
    = \mathcal{A} \otimes (\mathcal{B} \otimes \mathcal{C})
  \]
  is an isomorphism. Between the two un-sheafified middle terms sits the
  sheafification of the presheaf associator
  \[
    \alpha^p_{a,b,c} :
    (a \otimes_p b) \otimes_p c
    \xrightarrow{\ \sim\ }
    a \otimes_p (b \otimes_p c)
  \]
  (\cref{lem:presheafModule_monoidal_mathlib}), an isomorphism because
  sheafification is a functor and \(\alpha^p\) is one. The associator is the
  composite
  \[
    \alpha_{\mathcal{A},\mathcal{B},\mathcal{C}}
    = (a \lhd \eta_{b \otimes_p c})^{\#}
      \circ (\alpha^p_{a,b,c})^{\#}
      \circ \bigl((\eta_{a \otimes_p b} \rhd c)^{\#}\bigr)^{-1},
  \]
  an isomorphism as a composite of isomorphisms. Naturality in
  \(\mathcal{A}, \mathcal{B}, \mathcal{C}\) follows from naturality of the unit,
  of the whiskerings, and of the presheaf associator.
\end{proof}

\subsection{Symmetric monoidal coherence by localization transfer}
\label{sec:sgr_localized_monoidal}

The associator (\cref{cor:sheafTensorObjAssoc}), braiding (\cref{def:tensorBraiding})
and unitors (\cref{def:tensorObjUnitIso}, \cref{def:tensorObjRightUnitor}) above are
the structural isomorphisms of a genuine \emph{symmetric monoidal} structure on
\(X.\Modules\). Rather than verify Mac\,Lane's coherence axioms --- pentagon,
triangle, hexagon, and \(\beta_{\mathbf{1},\mathbf{1}} = \mathrm{id}\) --- by hand
for these hand-built isomorphisms, we obtain them all at one stroke from the
following observation: \(X.\Modules\) is the \emph{localization} of
\(\PshMod(\mathcal{O}_X)\) at the class \(W\) of morphisms inverted by
sheafification, and a localization of a symmetric monoidal category at a
monoidal class of morphisms is again symmetric monoidal, with every coherence law
proved once and for all. The next blocks record the Mathlib machinery that supplies
this, the single new precondition the project must discharge, and the bridge
identifying the hand-built isomorphisms with the transferred ones.

\begin{theorem}[Monoidal structure on a monoidal localization]
  \label{thm:localizedMonoidal_mathlib}
  \lean{CategoryTheory.LocalizedMonoidal}
  \mathlibok
  \textit{Provided by Mathlib
  (\texttt{Mathlib.CategoryTheory.Localization.Monoidal.Basic},
  \texttt{.Braided}, \texttt{.Functor}).}
  Let \(\mathcal{D}\) be a monoidal category, \(W\) a \emph{monoidal} class of
  morphisms (\(W\) is
  multiplicative and stable under whiskering by an arbitrary object on the left
  and on the right), \(L : \mathcal{D} \to \mathcal{E}\) a localization of
  \(\mathcal{D}\) at \(W\), and \(\varepsilon : L(\mathbf{1}_{\mathcal{D}}) \cong u\)
  a choice of unit comparison. Then the type synonym
  \(\mathrm{LocalizedMonoidal}\,L\,W\,\varepsilon\) of \(\mathcal{E}\) carries a
  monoidal category structure --- tensor product, associator
  \(\alpha^{\mathrm{loc}}\), left and right unitors
  \(\lambda^{\mathrm{loc}}, \rho^{\mathrm{loc}}\) --- satisfying the pentagon and
  triangle axioms, and \(L\) becomes a strong monoidal functor onto it,
  with comparison isomorphism
  \(\mu_{X,Y} : L(X) \otimes L(Y) \xrightarrow{\sim} L(X \otimes Y)\) and
  componentwise formulas for the associator, the unitors, and the naturality of
  \(\mu\). If \(\mathcal{D}\) is braided (resp. symmetric), the
  synonym is braided (resp. symmetric) with braiding \(\beta^{\mathrm{loc}}\),
  satisfying the hexagon identities. Crucially the
  structure is carried by the type synonym, \emph{not} by \(\mathcal{E}\) itself,
  so it never collides with any pre-existing structure on \(\mathcal{E}\).
\end{theorem}

\begin{lemma}[Pentagon axiom]
  \label{lem:monoidal_pentagon_mathlib}
  \lean{CategoryTheory.MonoidalCategory.pentagon}
  \mathlibok
  \textit{Provided by Mathlib (\texttt{Mathlib.CategoryTheory.Monoidal.Category}).}
  In any monoidal category the associator satisfies the pentagon identity: for
  objects \(W, X, Y, Z\) the two ways of reassociating
  \(((W \otimes X) \otimes Y) \otimes Z\) into
  \(W \otimes (X \otimes (Y \otimes Z))\) agree,
  \[
    (\alpha_{W,X,Y} \rhd Z) \,;\, \alpha_{W, X \otimes Y, Z}
      \,;\, (W \lhd \alpha_{X,Y,Z})
      = \alpha_{W \otimes X, Y, Z} \,;\, \alpha_{W, X, Y \otimes Z}.
  \]
\end{lemma}

\begin{lemma}[Triangle axiom]
  \label{lem:monoidal_triangle_mathlib}
  \lean{CategoryTheory.MonoidalCategory.triangle}
  \mathlibok
  \textit{Provided by Mathlib (\texttt{Mathlib.CategoryTheory.Monoidal.Category}).}
  In any monoidal category, for objects \(X, Y\) the associator and the unitors
  satisfy
  \[
    \alpha_{X, \mathbf{1}, Y} \,;\, (X \lhd \lambda_Y) = \rho_X \rhd Y,
  \]
  relating the right unitor on the left factor to the left unitor on the right
  factor through the associator.
\end{lemma}

\begin{lemma}[Hexagon axiom]
  \label{lem:braided_hexagon_mathlib}
  \lean{CategoryTheory.BraidedCategory.hexagon_forward}
  \mathlibok
  \textit{Provided by Mathlib
  (\texttt{Mathlib.CategoryTheory.Monoidal.Braided.Basic}).}
  In any braided monoidal category the braiding and the associator satisfy the
  hexagon identity: for objects \(X, Y, Z\),
  \[
    \alpha_{X,Y,Z} \,;\, \beta_{X, Y \otimes Z} \,;\, \alpha_{Y,Z,X}
      = (\beta_{X,Y} \rhd Z) \,;\, \alpha_{Y,X,Z}
        \,;\, (Y \lhd \beta_{X,Z}).
  \]
\end{lemma}

\begin{lemma}[Braiding of the unit on the left]
  \label{lem:braiding_tensorUnit_left_mathlib}
  \lean{CategoryTheory.braiding_tensorUnit_left}
  \mathlibok
  \textit{Provided by Mathlib
  (\texttt{Mathlib.CategoryTheory.Monoidal.Braided.Basic}).}
  In a braided monoidal category, for every object \(X\),
  \(\beta_{\mathbf{1}, X} = \lambda_X \,;\, \rho_X^{-1}\): the braiding of the unit
  on the left is the composite of the left unitor with the inverse right unitor.
\end{lemma}

\begin{lemma}[Unitors agree on the unit]
  \label{lem:unitors_equal_mathlib}
  \lean{CategoryTheory.MonoidalCategory.unitors_equal}
  \mathlibok
  \textit{Provided by Mathlib
  (\texttt{Mathlib.CategoryTheory.Monoidal.CoherenceLemmas}).}
  In any monoidal category the two unitors of the unit object coincide,
  \(\lambda_{\mathbf{1}} = \rho_{\mathbf{1}}\). Together with
  \cref{lem:braiding_tensorUnit_left_mathlib} (at \(X = \mathbf{1}\)) this gives
  \(\beta_{\mathbf{1},\mathbf{1}} = \mathrm{id}\), the symmetry of the unit with
  itself used in the right-unit and commutativity reductions.
\end{lemma}

\begin{lemma}[Sheafification is a localization at the sheafification class]
  \label{lem:sheafification_isLocalization_mathlib}
  \lean{PresheafOfModules.inverseImage_W_toPresheaf_eq_inverseImage_isomorphisms}
  \mathlibok
  \textit{Provided by Mathlib
  (\texttt{Mathlib.Algebra.Category.ModuleCat.Sheaf.Localization}).}
  Module sheafification \((-)^{\#} : \PshMod(R_0) \to \ShMod(R)\) is a
  localization functor for the class \(W\) of morphisms of presheaves of modules
  whose underlying abelian-presheaf morphism is a local isomorphism for the
  Grothendieck topology: Mathlib identifies \(W\) with the inverse image of the
  isomorphism class under sheafification, so sheafification is a localization at
  \(W\). For \(X.\Modules\) this specialises (\cref{def:schemeModuleSheafification},
  \cref{lem:isIso_sheafification_map_iff}) to \(L = \) module sheafification with
  the same class \(W\). This is the localization hypothesis of
  \cref{thm:localizedMonoidal_mathlib}.
\end{lemma}

The one new precondition \cref{thm:localizedMonoidal_mathlib} requires is that the
sheafification class \(W\) be monoidal. Its substance is already in hand: stability
under right whiskering is the proved \cref{lem:snap_ztensor_whisker_localIso}.

\begin{definition}\leanok
[The sheafification class is monoidal]
  \label{def:W_isMonoidal}
  \lean{AlgebraicGeometry.Scheme.Modules.W_isMonoidal}
  \uses{lem:snap_ztensor_whisker_localIso, def:tensorBraiding,
    lem:presheafModule_monoidal_mathlib, lem:sheafification_isLocalization_mathlib}
  The class \(W\) of morphisms of presheaves of \(\mathcal{O}_X\)-modules inverted
  by sheafification is \emph{monoidal}: it is multiplicative and stable under
  whiskering by an arbitrary presheaf of modules on either side.
  Multiplicativity (containing identities, closed under composition) is inherited
  from the local-isomorphism class through the inverse-image construction. Stability under right
  whiskering --- if \(f \in W\) then \(f \rhd \mathcal{P} \in W\) --- is exactly
  \cref{lem:snap_ztensor_whisker_localIso}. Stability under left
  whiskering follows by conjugating with the symmetric braiding of the presheaf
  monoidal structure (\cref{lem:presheafModule_monoidal_mathlib}): since
  \(\mathcal{P} \lhd f\) equals \(\beta \,;\, (f \rhd \mathcal{P}) \,;\, \beta\) up
  to the braiding isomorphisms (which lie in \(W\)), right-whiskering stability
  transports to the left. By the definition of a monoidal class of morphisms it suffices
  to verify multiplicativity together with these two whiskering stabilities, which
  the above supply.
\end{definition}

\begin{definition}\leanok
[Unit comparison for the localized monoidal structure]
  \label{def:localizedMonoidalUnitIso}
  \lean{AlgebraicGeometry.Scheme.Modules.localizedMonoidalUnitIso}
  \uses{def:schemeModuleSheafification, def:unitModule, lem:moduleUnit_mathlib,
    lem:presheafModule_monoidal_mathlib}
  The unit comparison isomorphism
  \[
    \varepsilon :
      (\mathbf{1}_{\PshMod})^{\#}
      \;\xrightarrow{\ \sim\ }\; \mathbf{1}_X
  \]
  identifying the sheafification (\cref{def:schemeModuleSheafification}) of the
  presheaf monoidal unit --- the structure presheaf viewed as a module over itself
  (\cref{lem:presheafModule_monoidal_mathlib}) --- with the unit module
  \(\mathbf{1}_X\) of \(X.\Modules\) (\cref{def:unitModule},
  \cref{lem:moduleUnit_mathlib}). It is the datum \(\varepsilon\) required to
  instantiate \cref{thm:localizedMonoidal_mathlib}.
\end{definition}

\begin{definition}\leanok
[The localized symmetric monoidal structure on the module sheaves]
  \label{def:modulesLocalizedMonoidal}
  \lean{AlgebraicGeometry.Scheme.Modules.modulesLocalizedMonoidal}
  \uses{thm:localizedMonoidal_mathlib, lem:sheafification_isLocalization_mathlib,
    def:W_isMonoidal, def:localizedMonoidalUnitIso,
    lem:presheafModule_monoidal_mathlib, lem:monoidal_pentagon_mathlib,
    lem:monoidal_triangle_mathlib, lem:braided_hexagon_mathlib,
    lem:braiding_tensorUnit_left_mathlib, lem:unitors_equal_mathlib}
  Instantiating \cref{thm:localizedMonoidal_mathlib} with
  \(\mathcal{D} = \PshMod(\mathcal{O}_X)\) (symmetric monoidal,
  \cref{lem:presheafModule_monoidal_mathlib}), \(L = \) module sheafification,
  \(W\) the sheafification class (monoidal by \cref{def:W_isMonoidal}, a
  localization by \cref{lem:sheafification_isLocalization_mathlib}), and
  \(\varepsilon\) the unit comparison \cref{def:localizedMonoidalUnitIso}, yields
  the type synonym \(\mathrm{LocalizedMonoidal}\,L\,W\,\varepsilon\), canonically
  identified with \(X.\Modules\) as a category, and carrying a \emph{symmetric monoidal}
  category structure: associator \(\alpha^{\mathrm{loc}}\), unitors
  \(\lambda^{\mathrm{loc}}, \rho^{\mathrm{loc}}\), braiding \(\beta^{\mathrm{loc}}\),
  with the pentagon, triangle and hexagon identities holding by
  \cref{lem:monoidal_pentagon_mathlib}, \cref{lem:monoidal_triangle_mathlib},
  \cref{lem:braided_hexagon_mathlib} and \(\beta^{\mathrm{loc}}_{\mathbf{1},
  \mathbf{1}} = \mathrm{id}\) by
  \cref{lem:braiding_tensorUnit_left_mathlib}, \cref{lem:unitors_equal_mathlib}.
  The structure lives on the type synonym, deliberately not on \(X.\Modules\)
  itself, so it introduces no coherence clash with the rest of the development.
\end{definition}

\providecommand{\lotimes}{\mathbin{\otimes_{\mathrm{loc}}}}

The bridge between \cref{def:modulesLocalizedMonoidal} and the hand-built
isomorphisms of \cref{cor:sheafTensorObjAssoc}, \cref{def:tensorBraiding},
\cref{def:tensorObjUnitIso}, \cref{def:tensorObjRightUnitor} must respect a
\emph{type distinction}. The synonym
\(\mathrm{LocalizedMonoidal}\,L\,W\,\varepsilon\) shares the underlying
\emph{category} of \(X.\Modules\), but its tensor product
\[
  \mathcal{F} \lotimes \mathcal{G}
    \;=\; \bigl((\mathrm{tensorBifunctor}\,L\,W\,\varepsilon)(\mathcal{F})\bigr)(\mathcal{G})
\]
is \emph{not} the hand-built tensor product
\(\mathcal{F} \otimes_{\mathcal{O}_X} \mathcal{G} =
(\mathcal{F}^{\flat} \otimes_p \mathcal{G}^{\flat})^{\#}\)
(\cref{def:sheafTensorObj}); here \(\mathcal{F}^{\flat}\) denotes the underlying
presheaf of modules of a sheaf \(\mathcal{F}\) and \((-)^{\#} = L\) the module
sheafification (\cref{def:schemeModuleSheafification}). The two tensor products
carry the same objects to \emph{different} objects of the category, so the
hand-built associator and the localized associator have different sources and
targets: a bare equality \(\alpha = \alpha^{\mathrm{loc}}\) between them is not a
well-formed statement. The comparison must therefore be threaded through the
object-identification isomorphism \cref{def:tensorObjLocalizedIso} on every tensor
slot.

We record first the Mathlib comparison datum, then the object identification, and
then the bridge equalities themselves.

\begin{lemma}[Strong-monoidality comparison of a monoidal localization]
  \label{lem:localizationMonoidal_mu_mathlib}
  \lean{CategoryTheory.Localization.Monoidal.μ}
  \mathlibok
  \textit{Provided by Mathlib
  (\texttt{Mathlib.CategoryTheory.Localization.Monoidal.Basic},
  \texttt{.Braided}).}
  In the setting of \cref{thm:localizedMonoidal_mathlib}, write
  \(L' : \mathcal{D} \to \mathrm{LocalizedMonoidal}\,L\,W\,\varepsilon\) for the
  localization functor onto the synonym (equal to \(L\) on objects and
  morphisms). For objects \(P, Q\) of \(\mathcal{D}\) there is a comparison
  isomorphism
  \[
    \mu_{P,Q} : L'(P) \lotimes L'(Q)
      \;\xrightarrow{\ \sim\ }\; L'(P \otimes_p Q),
  \]
  natural in \(P\) and in \(Q\): for \(f : P_1 \to P_2\) and \(g : Q_1 \to Q_2\)
  \[
    \bigl(L'(f) \rhd L'(Q)\bigr) \,;\, \mu_{P_2,Q}
      = \mu_{P_1,Q} \,;\, L'(f \rhd Q),
    \qquad
    \bigl(L'(P) \lhd L'(g)\bigr) \,;\, \mu_{P,Q_2}
      = \mu_{P,Q_1} \,;\, L'(P \lhd g)
  \]
  (\texttt{μ\_natural\_left}, \texttt{μ\_natural\_right}). Through \(\mu\), the
  unit comparison \(\varepsilon' = \varepsilon\), and \(L'\) applied to the
  structure maps of \(\mathcal{D}\), the structure isomorphisms of the synonym
  obey the componentwise formulas
  \[
    (\alpha^{\mathrm{loc}}_{L'P_1, L'P_2, L'P_3})^{\phantom{1}}
      = (\mu_{P_1,P_2} \rhd \cdot)\,;\, \mu_{P_1 \otimes P_2, P_3}
        \,;\, L'(\alpha^p_{P_1,P_2,P_3})
        \,;\, \mu_{P_1, P_2 \otimes P_3}^{-1}
        \,;\, (\cdot \lhd \mu_{P_2,P_3}^{-1}),
  \]
  \[
    \lambda^{\mathrm{loc}}_{L'Q}
      = (\varepsilon'^{-1} \rhd L'Q)\,;\, \mu_{\mathbf{1},Q}
        \,;\, L'(\lambda^p_Q),
    \qquad
    \rho^{\mathrm{loc}}_{L'P}
      = (L'P \lhd \varepsilon'^{-1})\,;\, \mu_{P,\mathbf{1}}
        \,;\, L'(\rho^p_P),
  \]
  \[
    \beta^{\mathrm{loc}}_{L'P, L'Q}
      = \mu_{P,Q} \,;\, L'(\beta^p_{P,Q}) \,;\, \mu_{Q,P}^{-1}
  \]
  (\texttt{associator\_hom\_app}, \texttt{leftUnitor\_hom\_app},
  \texttt{rightUnitor\_hom\_app}, \texttt{braidingNatIso\_hom\_app}; in the last,
  \(\mu\) is the lax-monoidal structure map of \(L'\) and \(\mu^{-1}\) its
  oplax counterpart).
\end{lemma}

\begin{definition}
[Object identification: hand-built versus localized tensor]
  \label{def:tensorObjLocalizedIso}
  \lean{AlgebraicGeometry.Scheme.Modules.tensorObjLocalizedIso}
  \uses{def:sheafTensorObj, def:modulesLocalizedMonoidal,
    def:schemeModuleSheafification, lem:localizationMonoidal_mu_mathlib,
    lem:presheafModule_sheafification_mathlib}
  For sheaves of \(\mathcal{O}_X\)-modules \(\mathcal{F}, \mathcal{G}\) the
  \emph{object-identification isomorphism}
  \[
    i_{\mathcal{F},\mathcal{G}} :
      \mathcal{F} \otimes_{\mathcal{O}_X} \mathcal{G}
      \;\xrightarrow{\ \sim\ }\;
      \mathcal{F} \lotimes \mathcal{G}
  \]
  identifies the hand-built tensor product (\cref{def:sheafTensorObj}) with the
  localized one (\cref{def:modulesLocalizedMonoidal}). It is the composite
  \[
    (\mathcal{F}^{\flat} \otimes_p \mathcal{G}^{\flat})^{\#}
      \;\xrightarrow{\ \mu_{\mathcal{F}^{\flat},\,\mathcal{G}^{\flat}}^{-1}\ }\;
      (\mathcal{F}^{\flat})^{\#} \lotimes (\mathcal{G}^{\flat})^{\#}
      \;\xrightarrow{\ c_{\mathcal{F}} \lotimes c_{\mathcal{G}}\ }\;
      \mathcal{F} \lotimes \mathcal{G},
  \]
  where \(\mu\) is the strong-monoidality comparison
  (\cref{lem:localizationMonoidal_mu_mathlib}) at the underlying presheaves
  \(\mathcal{F}^{\flat}, \mathcal{G}^{\flat}\), and
  \(c_{\mathcal{F}} : (\mathcal{F}^{\flat})^{\#} \xrightarrow{\sim} \mathcal{F}\)
  is the reflective sheafification counit isomorphism
  (\cref{def:schemeModuleSheafification}) that identifies the sheafification of the
  underlying presheaf of a sheaf of modules with the sheaf itself. (Here
  \(c_{\mathcal{F}} \lotimes c_{\mathcal{G}}\) is the tensor of the two counit
  isomorphisms in the localized monoidal category.) It is an isomorphism as a
  composite of isomorphisms, and natural in \(\mathcal{F}, \mathcal{G}\).
\end{definition}

Two routes now connect the localized structure to the hand-built isomorphisms.
\textbf{Option~A}: \emph{redefine} the hand-built tensor product to be the
localized one, \(\mathcal{F} \otimes_{\mathcal{O}_X} \mathcal{G} :=
\mathcal{F} \lotimes \mathcal{G}\), and likewise take the structural isomorphisms
to be the synonym's, so that the coherence laws of
\cref{def:modulesLocalizedMonoidal} apply to them literally. This route is
\emph{rejected}: the hand-built \(\otimes_{\mathcal{O}_X}\) is downstream
load-bearing --- \cref{def:sheafTensorPow} (\(\mathrm{tensorPow}\)),
\cref{def:sheafModuleTwist} (\(\mathrm{moduleTensorPow}\)) and the section
multiplication \cref{def:sectionMul} all rely on the definitional equality
\((\mathcal{F}^{\flat} \otimes_p \mathcal{G}^{\flat})(\top) =
\mathcal{F}^{\flat}(\top) \otimes \mathcal{G}^{\flat}(\top)\) of the value on the
top open, so redefining \(\otimes_{\mathcal{O}_X}\) has a large definitional blast
radius. \textbf{Option~B} (chosen): keep the hand-built definitions and prove,
for each structural isomorphism, the well-typed \emph{commuting square} expressing
the hand-built map as the corresponding map of the synonym \emph{conjugated by the
object identifications} \(i_{-,-}\) of \cref{def:tensorObjLocalizedIso} on every
tensor slot. Each square reduces, via the componentwise formulas of
\cref{lem:localizationMonoidal_mu_mathlib} and the naturality of \(\mu\), to the
sheafified presheaf-level coherence; the localized coherence laws then transfer
across the identifications. The four bridge lemmas below are exactly these
squares.

\begin{lemma}[The hand-built associator is the localized associator]
  \label{lem:tensorObjAssoc_eq_localizedAssociator}
  \lean{AlgebraicGeometry.Scheme.Modules.tensorObjAssoc_eq_localizedAssociator}
  \uses{cor:sheafTensorObjAssoc, def:modulesLocalizedMonoidal,
    def:tensorObjLocalizedIso, lem:localizationMonoidal_mu_mathlib,
    lem:isIso_sheafification_whiskerRight_unit, thm:localizedMonoidal_mathlib}
  For sheaves of \(\mathcal{O}_X\)-modules \(\mathcal{A}, \mathcal{B},
  \mathcal{C}\), the hand-built associator
  \(\alpha_{\mathcal{A},\mathcal{B},\mathcal{C}}\) (\cref{cor:sheafTensorObjAssoc})
  equals the localized associator
  \(\alpha^{\mathrm{loc}}_{\mathcal{A},\mathcal{B},\mathcal{C}}\)
  (\cref{def:modulesLocalizedMonoidal}) conjugated by the object identifications
  \(i_{-,-}\) of \cref{def:tensorObjLocalizedIso} on all four tensor slots. Writing
  \[
    \Phi^{L}_{\mathcal{A},\mathcal{B},\mathcal{C}}
      = i_{\mathcal{A} \otimes_{\mathcal{O}_X} \mathcal{B},\,\mathcal{C}}
        \,;\, \bigl(i_{\mathcal{A},\mathcal{B}} \rhd \mathcal{C}\bigr)
      : (\mathcal{A} \otimes_{\mathcal{O}_X} \mathcal{B})
          \otimes_{\mathcal{O}_X} \mathcal{C}
        \xrightarrow{\sim}
        (\mathcal{A} \lotimes \mathcal{B}) \lotimes \mathcal{C},
  \]
  \[
    \Phi^{R}_{\mathcal{A},\mathcal{B},\mathcal{C}}
      = i_{\mathcal{A},\,\mathcal{B} \otimes_{\mathcal{O}_X} \mathcal{C}}
        \,;\, \bigl(\mathcal{A} \lhd i_{\mathcal{B},\mathcal{C}}\bigr)
      : \mathcal{A} \otimes_{\mathcal{O}_X}
          (\mathcal{B} \otimes_{\mathcal{O}_X} \mathcal{C})
        \xrightarrow{\sim}
        \mathcal{A} \lotimes (\mathcal{B} \lotimes \mathcal{C}),
  \]
  the bridge is the commuting square
  \[
    \Phi^{L}_{\mathcal{A},\mathcal{B},\mathcal{C}} \,;\,
      \alpha^{\mathrm{loc}}_{\mathcal{A},\mathcal{B},\mathcal{C}}
      \;=\;
      \alpha_{\mathcal{A},\mathcal{B},\mathcal{C}} \,;\,
      \Phi^{R}_{\mathcal{A},\mathcal{B},\mathcal{C}},
    \qquad\text{equivalently}\qquad
    \alpha_{\mathcal{A},\mathcal{B},\mathcal{C}}
      = \Phi^{L}_{\mathcal{A},\mathcal{B},\mathcal{C}} \,;\,
        \alpha^{\mathrm{loc}}_{\mathcal{A},\mathcal{B},\mathcal{C}} \,;\,
        \bigl(\Phi^{R}_{\mathcal{A},\mathcal{B},\mathcal{C}}\bigr)^{-1},
  \]
  a well-typed equation between isomorphisms
  \((\mathcal{A} \otimes_{\mathcal{O}_X} \mathcal{B}) \otimes_{\mathcal{O}_X}
  \mathcal{C} \to \mathcal{A} \otimes_{\mathcal{O}_X}
  (\mathcal{B} \otimes_{\mathcal{O}_X} \mathcal{C})\).
\end{lemma}

\begin{proof}
  \uses{cor:sheafTensorObjAssoc, def:modulesLocalizedMonoidal,
    def:tensorObjLocalizedIso, lem:localizationMonoidal_mu_mathlib,
    lem:isIso_sheafification_whiskerRight_unit, thm:localizedMonoidal_mathlib}
  By naturality of \(\alpha^{\mathrm{loc}}\) and of \(\mu\) along the counit
  isomorphisms \(c_{\mathcal{A}}, c_{\mathcal{B}}, c_{\mathcal{C}}\), it suffices to
  prove the square at objects of the form
  \(\mathcal{A} = (P)^{\#}, \mathcal{B} = (Q)^{\#}, \mathcal{C} = (R)^{\#}\) with
  \(P = \mathcal{A}^{\flat}\) etc.; the object identifications then absorb the
  counits, and the displayed square reduces to the Mathlib component formula
  \texttt{associator\_hom\_app} (\cref{lem:localizationMonoidal_mu_mathlib}):
  \(\alpha^{\mathrm{loc}}_{L'P, L'Q, L'R}\) is
  \((\mu_{P,Q} \rhd \cdot)\,;\, \mu_{P \otimes Q, R}\,;\, L'(\alpha^p_{P,Q,R})\,;\,
  \mu_{P, Q \otimes R}^{-1}\,;\,(\cdot \lhd \mu_{Q,R}^{-1})\). Expanding
  \(\Phi^{L}, \Phi^{R}\) through \(\mu^{-1}\) (the definition of \(i_{-,-}\),
  \cref{def:tensorObjLocalizedIso}) and cancelling the \(\mu\)'s against the
  \(\mu^{-1}\)'s by \texttt{μ\_natural\_left}/\texttt{μ\_natural\_right}, both sides
  collapse to the single composite carried by \(L'(\alpha^p_{P,Q,R}) =
  (\alpha^p_{P,Q,R})^{\#}\). The hand-built associator
  (\cref{cor:sheafTensorObjAssoc}) is by construction \((\alpha^p)^{\#}\)
  conjugated by the whiskered-unit comparison isomorphisms
  (\cref{lem:isIso_sheafification_whiskerRight_unit}), and those comparisons are
  exactly the components of \(\mu\); the two composites therefore coincide, which
  is the square.
\end{proof}

\begin{lemma}[The hand-built braiding is the localized braiding]
  \label{lem:tensorBraiding_eq_localizedBraiding}
  \lean{AlgebraicGeometry.Scheme.Modules.tensorBraiding_eq_localizedBraiding}
  \uses{def:tensorBraiding, def:modulesLocalizedMonoidal,
    def:tensorObjLocalizedIso, lem:localizationMonoidal_mu_mathlib,
    lem:presheafModule_monoidal_mathlib, thm:localizedMonoidal_mathlib}
  For sheaves of \(\mathcal{O}_X\)-modules \(\mathcal{F}, \mathcal{G}\), the
  hand-built braiding \(\beta_{\mathcal{F},\mathcal{G}}\)
  (\cref{def:tensorBraiding}) equals the localized braiding
  \(\beta^{\mathrm{loc}}_{\mathcal{F},\mathcal{G}}\)
  (\cref{def:modulesLocalizedMonoidal}) conjugated by the object identifications
  \(i_{-,-}\) of \cref{def:tensorObjLocalizedIso} on its two tensor slots: the
  commuting square
  \[
    i_{\mathcal{F},\mathcal{G}} \,;\,
      \beta^{\mathrm{loc}}_{\mathcal{F},\mathcal{G}}
      \;=\;
      \beta_{\mathcal{F},\mathcal{G}} \,;\, i_{\mathcal{G},\mathcal{F}},
    \qquad\text{equivalently}\qquad
    \beta_{\mathcal{F},\mathcal{G}}
      = i_{\mathcal{F},\mathcal{G}} \,;\,
        \beta^{\mathrm{loc}}_{\mathcal{F},\mathcal{G}} \,;\,
        i_{\mathcal{G},\mathcal{F}}^{-1},
  \]
  a well-typed equation between isomorphisms
  \(\mathcal{F} \otimes_{\mathcal{O}_X} \mathcal{G} \to
  \mathcal{G} \otimes_{\mathcal{O}_X} \mathcal{F}\).
\end{lemma}

\begin{proof}
  \uses{def:tensorBraiding, def:modulesLocalizedMonoidal,
    def:tensorObjLocalizedIso, lem:localizationMonoidal_mu_mathlib,
    lem:presheafModule_monoidal_mathlib, thm:localizedMonoidal_mathlib}
  By naturality of \(\beta^{\mathrm{loc}}\) along the counit isomorphisms
  \(c_{\mathcal{F}}, c_{\mathcal{G}}\) it suffices to prove the square at
  \(\mathcal{F} = (P)^{\#}, \mathcal{G} = (Q)^{\#}\) with \(P =
  \mathcal{F}^{\flat}, Q = \mathcal{G}^{\flat}\). There the Mathlib component
  formula \texttt{braidingNatIso\_hom\_app}
  (\cref{lem:localizationMonoidal_mu_mathlib}) gives
  \(\beta^{\mathrm{loc}}_{L'P, L'Q} = \mu_{P,Q}\,;\, L'(\beta^p_{P,Q})\,;\,
  \mu_{Q,P}^{-1}\), where \(\beta^p\) is the presheaf braiding
  (\cref{lem:presheafModule_monoidal_mathlib}). Expanding
  \(i_{\mathcal{F},\mathcal{G}}\) and \(i_{\mathcal{G},\mathcal{F}}\) through
  \(\mu^{-1}\) (\cref{def:tensorObjLocalizedIso}), the \(\mu_{P,Q}\) and
  \(\mu_{Q,P}^{-1}\) cancel against the comparison factors of the identifications,
  leaving \(L'(\beta^p_{P,Q}) = (\beta^p_{P,Q})^{\#}\). That sheafified presheaf
  braiding is, by definition, the hand-built braiding
  \(\beta_{\mathcal{F},\mathcal{G}}\) (\cref{def:tensorBraiding}); the square
  follows.
\end{proof}

\begin{lemma}[The hand-built left unitor is the localized left unitor]
  \label{lem:tensorObjUnitor_eq_localized}
  \lean{AlgebraicGeometry.Scheme.Modules.tensorObjUnitIso_eq_localizedLeftUnitor}
  \uses{def:tensorObjUnitIso, def:modulesLocalizedMonoidal,
    def:tensorObjLocalizedIso, def:localizedMonoidalUnitIso,
    lem:localizationMonoidal_mu_mathlib, thm:localizedMonoidal_mathlib}
  For a sheaf of \(\mathcal{O}_X\)-modules \(\mathcal{G}\), the hand-built left
  unitor \(\lambda_{\mathcal{G}}\) (\cref{def:tensorObjUnitIso}) equals the
  localized left unitor \(\lambda^{\mathrm{loc}}_{\mathcal{G}}\)
  (\cref{def:modulesLocalizedMonoidal}) precomposed with the object identification
  \(i_{\mathbf{1}_X,\mathcal{G}}\) of \cref{def:tensorObjLocalizedIso}. The
  localized monoidal unit is the unit module
  \(\mathbf{1}_X = \mathbf{1}_{\mathrm{loc}}\) (the codomain of \(\varepsilon\),
  \cref{def:localizedMonoidalUnitIso}), so
  \(i_{\mathbf{1}_X,\mathcal{G}} : \mathbf{1}_X \otimes_{\mathcal{O}_X}
  \mathcal{G} \xrightarrow{\sim} \mathbf{1}_{\mathrm{loc}} \lotimes \mathcal{G}\),
  and the bridge is the well-typed identity
  \[
    \lambda_{\mathcal{G}}
      \;=\; i_{\mathbf{1}_X,\mathcal{G}} \,;\,
        \lambda^{\mathrm{loc}}_{\mathcal{G}}
    \;:\; \mathbf{1}_X \otimes_{\mathcal{O}_X} \mathcal{G}
      \xrightarrow{\sim} \mathcal{G}
  \]
  (the target slot is \(\mathcal{G}\) on both sides, so it carries the identity).
\end{lemma}

\begin{proof}
  \uses{def:tensorObjUnitIso, def:modulesLocalizedMonoidal,
    def:tensorObjLocalizedIso, def:localizedMonoidalUnitIso,
    lem:localizationMonoidal_mu_mathlib, thm:localizedMonoidal_mathlib}
  By naturality of \(\lambda^{\mathrm{loc}}\) along the counit \(c_{\mathcal{G}}\)
  it suffices to treat \(\mathcal{G} = (Q)^{\#}\) with \(Q = \mathcal{G}^{\flat}\).
  The Mathlib component formula \texttt{leftUnitor\_hom\_app}
  (\cref{lem:localizationMonoidal_mu_mathlib}) gives
  \(\lambda^{\mathrm{loc}}_{L'Q} = (\varepsilon'^{-1} \rhd L'Q)\,;\,
  \mu_{\mathbf{1},Q}\,;\, L'(\lambda^p_Q)\), where \(\varepsilon' = \varepsilon\)
  is the unit comparison (\cref{def:localizedMonoidalUnitIso}) and \(\lambda^p\)
  the presheaf left unitor. Expanding \(i_{\mathbf{1}_X,\mathcal{G}}\) through
  \(\mu_{\mathbf{1},Q}^{-1}\) and the counit on the unit slot
  (\cref{def:tensorObjLocalizedIso}), the comparison \(\mu_{\mathbf{1},Q}\) cancels
  and the \(\varepsilon'\) factors absorb the unit-object counit, leaving
  \(L'(\lambda^p_Q) = (\lambda^p_Q)^{\#}\) followed by the reflective counit. That
  is exactly the hand-built left unitor (\cref{def:tensorObjUnitIso}), which is the
  sheafified presheaf left unitor composed with the counit. The identity follows.
\end{proof}

\begin{lemma}[The hand-built right unitor is the localized right unitor]
  \label{lem:tensorObjRightUnitor_eq_localized}
  \lean{AlgebraicGeometry.Scheme.Modules.tensorObjRightUnitor_eq_localizedRightUnitor}
  \uses{def:tensorObjRightUnitor, def:modulesLocalizedMonoidal,
    def:tensorObjLocalizedIso, lem:tensorBraiding_eq_localizedBraiding,
    lem:tensorObjUnitor_eq_localized, lem:localizationMonoidal_mu_mathlib,
    thm:localizedMonoidal_mathlib}
  For a sheaf of \(\mathcal{O}_X\)-modules \(\mathcal{G}\), the hand-built right
  unitor \(\rho_{\mathcal{G}}\) (\cref{def:tensorObjRightUnitor}) equals the
  localized right unitor \(\rho^{\mathrm{loc}}_{\mathcal{G}}\)
  (\cref{def:modulesLocalizedMonoidal}) precomposed with the object identification
  \(i_{\mathcal{G},\mathbf{1}_X}\) of \cref{def:tensorObjLocalizedIso}. With
  \(\mathbf{1}_X = \mathbf{1}_{\mathrm{loc}}\)
  (\cref{def:localizedMonoidalUnitIso}), so
  \(i_{\mathcal{G},\mathbf{1}_X} : \mathcal{G} \otimes_{\mathcal{O}_X}
  \mathbf{1}_X \xrightarrow{\sim} \mathcal{G} \lotimes \mathbf{1}_{\mathrm{loc}}\),
  the bridge is the well-typed identity
  \[
    \rho_{\mathcal{G}}
      \;=\; i_{\mathcal{G},\mathbf{1}_X} \,;\,
        \rho^{\mathrm{loc}}_{\mathcal{G}}
    \;:\; \mathcal{G} \otimes_{\mathcal{O}_X} \mathbf{1}_X
      \xrightarrow{\sim} \mathcal{G}
  \]
  (the target slot is \(\mathcal{G}\) on both sides, carrying the identity).
\end{lemma}

\begin{proof}
  \uses{def:tensorObjRightUnitor, def:modulesLocalizedMonoidal,
    def:tensorObjLocalizedIso, lem:tensorBraiding_eq_localizedBraiding,
    lem:tensorObjUnitor_eq_localized, lem:braiding_tensorUnit_left_mathlib,
    lem:localizationMonoidal_mu_mathlib}
  The hand-built right unitor is \(\rho_{\mathcal{G}} = \lambda_{\mathcal{G}}
  \circ \beta_{\mathcal{G},\mathbf{1}_X}\) (\cref{def:tensorObjRightUnitor},
  diagrammatically \(\rho_{\mathcal{G}} = \beta_{\mathcal{G},\mathbf{1}_X}\,;\,
  \lambda_{\mathcal{G}}\)). In the symmetric localized structure the right unitor
  likewise satisfies \(\rho^{\mathrm{loc}}_{\mathcal{G}} =
  \beta^{\mathrm{loc}}_{\mathcal{G},\mathbf{1}_{\mathrm{loc}}}\,;\,
  \lambda^{\mathrm{loc}}_{\mathcal{G}}\) (the symmetric-category relation
  \cref{lem:braiding_tensorUnit_left_mathlib} at the unit, equivalently the
  componentwise formula \texttt{rightUnitor\_hom\_app} of
  \cref{lem:localizationMonoidal_mu_mathlib}). Substitute the braiding bridge
  \(\beta_{\mathcal{G},\mathbf{1}_X} = i_{\mathcal{G},\mathbf{1}_X}\,;\,
  \beta^{\mathrm{loc}}_{\mathcal{G},\mathbf{1}_{\mathrm{loc}}}\,;\,
  i_{\mathbf{1}_X,\mathcal{G}}^{-1}\)
  (\cref{lem:tensorBraiding_eq_localizedBraiding}) and the left-unitor bridge
  \(\lambda_{\mathcal{G}} = i_{\mathbf{1}_X,\mathcal{G}}\,;\,
  \lambda^{\mathrm{loc}}_{\mathcal{G}}\)
  (\cref{lem:tensorObjUnitor_eq_localized}) into
  \(\rho_{\mathcal{G}} = \beta_{\mathcal{G},\mathbf{1}_X}\,;\,
  \lambda_{\mathcal{G}}\); the inner
  \(i_{\mathbf{1}_X,\mathcal{G}}^{-1}\,;\, i_{\mathbf{1}_X,\mathcal{G}}\) cancels,
  leaving \(\rho_{\mathcal{G}} = i_{\mathcal{G},\mathbf{1}_X}\,;\,
  \beta^{\mathrm{loc}}_{\mathcal{G},\mathbf{1}_{\mathrm{loc}}}\,;\,
  \lambda^{\mathrm{loc}}_{\mathcal{G}} = i_{\mathcal{G},\mathbf{1}_X}\,;\,
  \rho^{\mathrm{loc}}_{\mathcal{G}}\), the claim.
\end{proof}

\begin{lemma}\leanok
[Tensor-power comparison isomorphism]
  \label{lem:sheafTensorPow_add}
  \lean{AlgebraicGeometry.Scheme.Modules.tensorPowAdd}
  \uses{def:sheafTensorPow, def:sheafTensorObj}
  % SOURCE: [Stacks Project], "Sheaves of Modules", \S "Invertible modules"
  % (Tag 01CU, following `definition-powers') (read from
  % references/stacks-modules.tex, L4249-L4254).
  % SOURCE QUOTE:
  % "With this definition we have canonical isomorphisms
  %  $\mathcal{L}^{\otimes n} \otimes_{\mathcal{O}_X}
  %  \mathcal{L}^{\otimes m} \to
  %  \mathcal{L}^{\otimes n + m}$, and these isomorphisms
  %  satisfy a commutativity and an associativity constraint
  %  (formulation omitted)."
  \textit{Source: [Stacks Project], "Sheaves of Modules", Tag 01CU (canonical
  power isomorphisms).}
  For a line bundle \(\mathcal{L}\) and \(m, m' \in \mathbb{N}\) there is a
  canonical isomorphism of sheaves of \(\mathcal{O}_X\)-modules
  \[
    \mu_{m,m'} :
    \mathcal{L}^{\otimes m} \otimes_{\mathcal{O}_X} \mathcal{L}^{\otimes m'}
    \;\xrightarrow{\ \sim\ }\; \mathcal{L}^{\otimes (m+m')},
  \]
  natural in \(\mathcal{L}\), and these isomorphisms satisfy a commutativity
  constraint (\(\mu_{m,m'}\) agrees with \(\mu_{m',m}\) after the braiding) and
  an associativity constraint (the two ways of combining
  \(\mathcal{L}^{\otimes m} \otimes \mathcal{L}^{\otimes m'} \otimes
  \mathcal{L}^{\otimes m''}\) into \(\mathcal{L}^{\otimes (m+m'+m'')}\) agree).
\end{lemma}

\begin{proof}
  \uses{def:sheafTensorPow, def:sheafTensorObj, cor:sheafTensorObjAssoc,
    def:tensorObjUnitIso, def:tensorBraiding,
    lem:isIso_sheafification_whiskerRight_unit,
    lem:presheafModule_monoidal_mathlib,
    lem:presheafModule_sheafification_mathlib}
    \leanok
  We do \emph{not} build a full monoidal structure on the whole category
  \(\ShMod(\mathcal{O}_X)\) --- none is constructed here --- but only assemble
  the comparison \(\mu_{m,m'}\) from those structural isomorphisms of
  \cref{def:sheafTensorObj} that the induction actually uses: the associator,
  the braiding, and the left unitor.

  \textbf{The structural maps.} The braiding and the unitors descend to
  \(\ShMod(\mathcal{O}_X)\) directly: each is the sheafification of the
  corresponding isomorphism of the presheaf symmetric monoidal structure
  (\cref{lem:presheafModule_monoidal_mathlib}), and sheafification, being a
  functor, carries isomorphisms to isomorphisms. The \emph{associator} is not
  the sheafification of a single presheaf isomorphism, because the iterated
  sheaf tensor products
  \((\mathcal{A} \otimes \mathcal{B}) \otimes \mathcal{C}\) and
  \(\mathcal{A} \otimes (\mathcal{B} \otimes \mathcal{C})\) each carry an inner
  sheafification (\cref{def:sheafTensorObj}); it is supplied instead by
  \cref{cor:sheafTensorObjAssoc}, whose construction clears those inner
  sheafifications using the strong-monoidality comparison
  \cref{lem:isIso_sheafification_whiskerRight_unit}. These are the only
  isomorphisms of sheaves of modules the construction invokes; right-exactness
  of the tensor product is never used.

  \textbf{Construction of \(\mu_{m,m'}\).} Fix \(m'\) and induct on \(m\). For
  \(m = 0\),
  \[
    \mathcal{L}^{\otimes 0} \otimes \mathcal{L}^{\otimes m'}
      = \mathcal{O}_X \otimes \mathcal{L}^{\otimes m'}
      \xrightarrow{\ \sim\ } \mathcal{L}^{\otimes m'}
  \]
  is the left unitor of \cref{def:sheafTensorObj}. For the inductive step the
  recursive definition
  \(\mathcal{L}^{\otimes (m+1)} = \mathcal{L}^{\otimes m} \otimes \mathcal{L}\)
  (\cref{def:sheafTensorPow}) gives
  \[
    \mathcal{L}^{\otimes (m+1)} \otimes \mathcal{L}^{\otimes m'}
    = (\mathcal{L}^{\otimes m} \otimes \mathcal{L}) \otimes
      \mathcal{L}^{\otimes m'}
    \xrightarrow{\ \alpha\ } \mathcal{L}^{\otimes m} \otimes
      (\mathcal{L} \otimes \mathcal{L}^{\otimes m'})
    \xrightarrow{\ \mathrm{id} \otimes \beta\ } \mathcal{L}^{\otimes m} \otimes
      (\mathcal{L}^{\otimes m'} \otimes \mathcal{L}),
  \]
  using the associator \(\alpha\) (\cref{cor:sheafTensorObjAssoc}) and then the
  braiding \(\beta\) of \cref{def:sheafTensorObj} whiskered into the left
  factor. The inverse associator \(\alpha^{-1}\)
  (\cref{cor:sheafTensorObjAssoc}) and the inductive comparison \(\mu_{m,m'}\)
  whiskered by \(\mathcal{L}\) on the right then carry this to
  \[
    (\mathcal{L}^{\otimes m} \otimes \mathcal{L}^{\otimes m'}) \otimes
      \mathcal{L}
    \xrightarrow{\ \mu_{m,m'} \rhd \mathcal{L}\ }
    \mathcal{L}^{\otimes (m + m')} \otimes \mathcal{L}
    = \mathcal{L}^{\otimes ((m + m') + 1)},
  \]
  the final equality being the recursive definition. The reindexing
  \((m + m') + 1 = (m + 1) + m'\) identifies the target with
  \(\mathcal{L}^{\otimes ((m+1) + m')}\), and the composite is
  \(\mu_{m+1,m'}\), an isomorphism as a composite of isomorphisms. Naturality in
  \(\mathcal{L}\) holds because each constituent comparison is natural.

  \textbf{Commutativity and associativity constraints.} The braiding and
  unitors are the sheafifications of the symmetric-monoidal coherence data of
  \(\PshMod(\mathcal{O}_X)\) (\cref{lem:presheafModule_monoidal_mathlib}), and
  the associator is intertwined with the presheaf associator by the comparison
  isomorphisms of \cref{lem:isIso_sheafification_whiskerRight_unit} through the
  construction of \cref{cor:sheafTensorObjAssoc}. Both the commutativity square
  (\(\mu_{m,m'}\) versus \(\mu_{m',m}\) postcomposed with the braiding) and the
  associativity pentagon (the two bracketings of
  \(\mathcal{L}^{\otimes m} \otimes \mathcal{L}^{\otimes m'} \otimes
  \mathcal{L}^{\otimes m''} \to \mathcal{L}^{\otimes (m+m'+m'')}\)) are therefore
  images, under the sheafification functor together with these comparison
  isomorphisms, of the corresponding coherence diagrams of the presheaf
  symmetric monoidal structure, which commute by Mac\,Lane coherence. Since a
  functor preserves commuting diagrams and the comparison isomorphisms are
  natural, the constraints hold for the family \(\{\mu_{m,m'}\}\). We emphasise
  that this argument uses only the single strong-monoidality comparison
  \cref{lem:isIso_sheafification_whiskerRight_unit} and the resulting associator
  \cref{cor:sheafTensorObjAssoc}, not a full monoidal structure on
  \(\ShMod(\mathcal{O}_X)\).
\end{proof}

The four coherence identities of the tensor-power comparison
\(\{\mu_{m,m'}\}\) are isolated below as named lemmas, each pinning one of the
recursive branches or symmetric-monoidal constraints used in the proof of
\cref{lem:sheafTensorPow_add}. They are the structural half of the graded-monoid
coherence \cref{lem:sectionMul_coherent}: the section-multiplication identities
factor through them after the section-level cores
(\cref{lem:sectionMul_unitor_core} and its analogues) discharge the
sheafification-unit bookkeeping. They mirror, branch for branch, the four unitality, associativity, and
commutativity identities of Mathlib's abstract tensor-power construction, with the
sheaf-level associator, braiding and unitors in place of the strict ones.

\begin{lemma}\leanok
[Left-unit branch of the tensor-power comparison]
  \label{lem:tensorPowAdd_zero}
  \lean{AlgebraicGeometry.Scheme.Modules.tensorPowAdd_zero}
  \uses{lem:sheafTensorPow_add, def:tensorObjUnitIso, def:sheafTensorPow}
  For \(n \in \mathbb{N}\) the comparison \(\mu_{0,n}\) of
  \cref{lem:sheafTensorPow_add} factors as the left unitor of
  \cref{def:tensorObjUnitIso} followed by the canonical index-equality
  isomorphism:
  \[
    \mu_{0,n}
      \;=\;
      \lambda_{\mathcal{L}^{\otimes n}}
      \;\textrm{followed by}\;
      \bigl(\mathcal{L}^{\otimes n}
        \xrightarrow{\ \sim\ } \mathcal{L}^{\otimes (0+n)}\bigr),
  \]
  the second arrow being the isomorphism induced by the index equality
  \(n = 0+n\) under the tensor-power functor (\cref{def:sheafTensorPow}).
\end{lemma}

\begin{proof}
  \uses{lem:sheafTensorPow_add}
  \leanok
  This is exactly the base branch of the recursion defining \(\mu_{m,n}\) on the
  first index \(m\) in \cref{lem:sheafTensorPow_add}. Since the first index is the
  literal \(0\), the recursion reduces to its base case, which is by construction
  the displayed composite; no unfolding of the successor branch is required.
\end{proof}

\begin{lemma}\leanok
[Right-unit branch of the tensor-power comparison]
  \label{lem:tensorPowAdd_rightUnit}
  \lean{AlgebraicGeometry.Scheme.Modules.tensorPowAdd_rightUnit}
  \uses{lem:sheafTensorPow_add, def:tensorObjRightUnitor, def:tensorObjUnitIso,
    def:tensorBraiding, lem:tensorPowAdd_zero, cor:sheafTensorObjAssoc,
    def:sheafTensorPow, lem:presheafModule_monoidal_mathlib}
  For \(n \in \mathbb{N}\) the comparison \(\mu_{n,0}\) of
  \cref{lem:sheafTensorPow_add} factors as the right unitor of
  \cref{def:tensorObjRightUnitor} followed by the canonical index-equality
  isomorphism:
  \[
    \mu_{n,0}
      \;=\;
      \rho_{\mathcal{L}^{\otimes n}}
      \;\textrm{followed by}\;
      \bigl(\mathcal{L}^{\otimes n}
        \xrightarrow{\ \sim\ } \mathcal{L}^{\otimes (n+0)}\bigr),
  \]
  the second arrow being induced by \(n = n+0\) (\cref{def:sheafTensorPow}).
  Unlike \cref{lem:tensorPowAdd_zero}, the first index here is the \emph{variable}
  \(n\), so the recursion does not reduce definitionally; the identity is proved by
  induction on \(n\) through the successor branch.
\end{lemma}

\begin{proof}
  \uses{lem:sheafTensorPow_add, def:tensorObjRightUnitor, def:tensorObjUnitIso,
    def:tensorBraiding, lem:tensorPowAdd_zero, cor:sheafTensorObjAssoc,
    def:modulesLocalizedMonoidal, lem:monoidal_triangle_mathlib,
    lem:braided_hexagon_mathlib, lem:tensorObjAssoc_eq_localizedAssociator,
    lem:tensorBraiding_eq_localizedBraiding, lem:tensorObjUnitor_eq_localized,
    lem:tensorObjRightUnitor_eq_localized}
  Induct on \(n\).

  \emph{Base \(n = 0\).} Here \(\mathcal{L}^{\otimes 0} = \mathbf{1}_X\). By
  \cref{lem:tensorPowAdd_zero}, \(\mu_{0,0}\) is
  \(\lambda_{\mathbf{1}_X}\) followed by the index isomorphism; the claim is that
  it equals \(\rho_{\mathbf{1}_X}\) followed by the same index isomorphism, i.e.\
  that \(\lambda_{\mathbf{1}_X} = \rho_{\mathbf{1}_X}\). By
  \cref{def:tensorObjRightUnitor},
  \(\rho_{\mathbf{1}_X} = \lambda_{\mathbf{1}_X} \circ
  \beta_{\mathbf{1}_X,\mathbf{1}_X}\); and the braiding of the unit with itself is
  the identity, because in the presheaf symmetric monoidal structure
  (\cref{lem:presheafModule_monoidal_mathlib}) the symmetry on the unit object is
  the identity and sheafification, being a functor, carries the identity to the
  identity (\cref{def:tensorBraiding}). Hence
  \(\rho_{\mathbf{1}_X} = \lambda_{\mathbf{1}_X}\).

  \emph{Step \(n = k+1\).} Write \(A = \mathcal{L}^{\otimes k}\) and
  \(\mathcal{L}\) for the generating bundle, so
  \(\mathcal{L}^{\otimes (k+1)} = A \otimes_{\mathcal{O}_X} \mathcal{L}\) and
  \(\mathcal{L}^{\otimes 0} = \mathbf{1}_X\). The successor branch of
  \cref{lem:sheafTensorPow_add} expresses \(\mu_{k+1,0}\), up to the index
  isomorphism \((k+0)+1 = (k+1)+0\), as the composite
  \[
    (A \otimes \mathcal{L}) \otimes \mathbf{1}_X
      \xrightarrow{\ \alpha\ }
      A \otimes (\mathcal{L} \otimes \mathbf{1}_X)
      \xrightarrow{\ A \lhd \beta_{\mathcal{L},\mathbf{1}_X}\ }
      A \otimes (\mathbf{1}_X \otimes \mathcal{L})
      \xrightarrow{\ \alpha^{-1}\ }
      (A \otimes \mathbf{1}_X) \otimes \mathcal{L}
      \xrightarrow{\ \mu_{k,0} \rhd \mathcal{L}\ }
      A \otimes \mathcal{L},
  \]
  with \(\alpha\) the associator (\cref{cor:sheafTensorObjAssoc}) and \(\beta\) the
  braiding (\cref{def:tensorBraiding}). By the inductive hypothesis
  \(\mu_{k,0} = \rho_A\) followed by the index isomorphism, so after cancelling
  the index isomorphisms the composite is
  \[
    (\rho_A \rhd \mathcal{L}) \circ \alpha^{-1}_{A,\mathbf{1}_X,\mathcal{L}}
      \circ (A \lhd \beta_{\mathcal{L},\mathbf{1}_X})
      \circ \alpha_{A,\mathcal{L},\mathbf{1}_X}.
  \]
  This is a formal composite of structural isomorphisms of a symmetric monoidal
  category, and it equals \(\rho_{A \otimes \mathcal{L}}\): the identity
  \[
    \rho_{A \otimes B}
      = (\rho_A \rhd B) \circ \alpha^{-1}_{A,\mathbf{1}_X,B}
        \circ (A \lhd \beta_{B,\mathbf{1}_X})
        \circ \alpha_{A,B,\mathbf{1}_X}
  \]
  (with \(B = \mathcal{L}\)) is a consequence of the triangle and hexagon axioms,
  hence holds in any symmetric monoidal category. The localized symmetric monoidal
  structure of \cref{def:modulesLocalizedMonoidal} \emph{is} such a structure on
  \(X.\Modules\), so the identity holds there directly by the triangle law
  \cref{lem:monoidal_triangle_mathlib} and the hexagon law
  \cref{lem:braided_hexagon_mathlib}, applied to the localized associator, braiding
  and unitors. The bridge lemmas \cref{lem:tensorObjAssoc_eq_localizedAssociator},
  \cref{lem:tensorBraiding_eq_localizedBraiding},
  \cref{lem:tensorObjUnitor_eq_localized} and
  \cref{lem:tensorObjRightUnitor_eq_localized} rewrite the hand-built
  \(\alpha, \beta, \lambda, \rho\) appearing in the composite as their localized
  counterparts, after which the displayed equality is the coherence law itself.
  Thus \(\mu_{k+1,0} = \rho_{\mathcal{L}^{\otimes (k+1)}}\) followed by the index
  isomorphism, closing the induction.
\end{proof}

\begin{lemma}\leanok
[Commutativity constraint of the tensor-power comparison]
  \label{lem:tensorPowAdd_braiding}
  \lean{AlgebraicGeometry.Scheme.Modules.tensorPowAdd_braiding}
  \uses{lem:sheafTensorPow_add, def:tensorBraiding, cor:sheafTensorObjAssoc,
    lem:tensorPowAdd_zero, lem:tensorPowAdd_rightUnit, def:sheafTensorPow,
    lem:presheafModule_monoidal_mathlib}
  For \(n_a, n_b \in \mathbb{N}\) the comparison \(\mu_{n_a,n_b}\) of
  \cref{lem:sheafTensorPow_add} agrees with \(\mu_{n_b,n_a}\) after braiding and
  reindexing:
  \[
    \mu_{n_a,n_b}
      \;=\;
      \beta_{\mathcal{L}^{\otimes n_a},\,\mathcal{L}^{\otimes n_b}}
      \;\textrm{followed by}\;
      \mu_{n_b,n_a}
      \;\textrm{followed by}\;
      \bigl(\mathcal{L}^{\otimes (n_b+n_a)}
        \xrightarrow{\ \sim\ } \mathcal{L}^{\otimes (n_a+n_b)}\bigr),
  \]
  the final arrow induced by \(n_b + n_a = n_a + n_b\)
  (\cref{def:sheafTensorPow}), with \(\beta\) the braiding
  (\cref{def:tensorBraiding}). This is the commutativity constraint asserted in
  \cref{lem:sheafTensorPow_add}.
\end{lemma}

\begin{proof}
  \uses{lem:sheafTensorPow_add, def:tensorBraiding, cor:sheafTensorObjAssoc,
    lem:tensorPowAdd_zero, lem:tensorPowAdd_rightUnit,
    lem:presheafModule_monoidal_mathlib, def:modulesLocalizedMonoidal,
    lem:braided_hexagon_mathlib, lem:tensorObjAssoc_eq_localizedAssociator,
    lem:tensorBraiding_eq_localizedBraiding}
  Induct on \(n_a\), keeping \(n_b\) fixed.

  \emph{Base \(n_a = 0\).} By \cref{lem:tensorPowAdd_zero}, \(\mu_{0,n_b}\) is the
  left unitor \(\lambda_{\mathcal{L}^{\otimes n_b}}\) followed by the index
  isomorphism. By \cref{lem:tensorPowAdd_rightUnit}, \(\mu_{n_b,0}\) is the right
  unitor \(\rho_{\mathcal{L}^{\otimes n_b}}\) followed by the index isomorphism, and
  by \cref{def:tensorObjRightUnitor} the right unitor is
  \(\lambda_{\mathcal{L}^{\otimes n_b}} \circ
  \beta_{\mathcal{L}^{\otimes n_b},\mathbf{1}_X}\). Thus the claimed right-hand
  side reduces, after cancelling index isomorphisms, to
  \[
    \lambda_{\mathcal{L}^{\otimes n_b}}
      \circ \beta_{\mathcal{L}^{\otimes n_b},\mathbf{1}_X}
      \circ \beta_{\mathbf{1}_X,\mathcal{L}^{\otimes n_b}}
      = \lambda_{\mathcal{L}^{\otimes n_b}},
  \]
  the braiding composed with its reverse being the identity by symmetry of the
  presheaf monoidal structure (\cref{lem:presheafModule_monoidal_mathlib}); this
  equals \(\mu_{0,n_b}\), proving the base case.

  \emph{Step \(n_a = k+1\).} Expanding \(\mu_{k+1,n_b}\) by the successor branch of
  \cref{lem:sheafTensorPow_add} and \(\mu_{n_b,k+1}\) on the other side, the two
  sides differ by a formal rearrangement of associators and braidings whose
  underlying diagram is one of the hexagon coherence diagrams for the symmetric
  braiding, together with the inductive hypothesis applied to \(\mu_{k,n_b}\)
  versus \(\mu_{n_b,k}\). This diagram commutes in any symmetric monoidal
  category; the localized structure of \cref{def:modulesLocalizedMonoidal} is one
  on \(X.\Modules\), so it commutes there by the hexagon law
  \cref{lem:braided_hexagon_mathlib} applied to the localized associator and
  braiding, the bridge lemmas
  \cref{lem:tensorObjAssoc_eq_localizedAssociator} and
  \cref{lem:tensorBraiding_eq_localizedBraiding} rewriting the hand-built
  \(\alpha\) and \(\beta\) as their localized counterparts. This yields the
  identity for \(n_a = k+1\) and closes the induction.
\end{proof}

\begin{lemma}\leanok
[Associativity constraint of the tensor-power comparison]
  \label{lem:tensorPowAdd_assoc}
  \lean{AlgebraicGeometry.Scheme.Modules.tensorPowAdd_assoc}
  \uses{lem:sheafTensorPow_add, cor:sheafTensorObjAssoc, def:sheafTensorPow,
    lem:presheafModule_monoidal_mathlib, lem:tensorPowAdd_zero}
  For \(n_a, n_b, n_c \in \mathbb{N}\) the two ways of combining
  \(\mathcal{L}^{\otimes n_a} \otimes \mathcal{L}^{\otimes n_b} \otimes
  \mathcal{L}^{\otimes n_c}\) into a single tensor power agree: the comparison
  isomorphisms of \cref{lem:sheafTensorPow_add} satisfy
  \[
    \bigl(\mathcal{L}^{\otimes ((n_a+n_b)+n_c)}
      \xrightarrow{\ \sim\ } \mathcal{L}^{\otimes (n_a+(n_b+n_c))}\bigr)
      \circ \mu_{n_a+n_b,\,n_c}
      \circ \bigl(\mu_{n_a,n_b} \rhd \mathcal{L}^{\otimes n_c}\bigr)
    \;=\;
    \mu_{n_a,\,n_b+n_c}
      \circ \bigl(\mathcal{L}^{\otimes n_a} \lhd \mu_{n_b,n_c}\bigr)
      \circ \alpha_{\mathcal{L}^{\otimes n_a},\,\mathcal{L}^{\otimes n_b},\,
        \mathcal{L}^{\otimes n_c}},
  \]
  where \(\alpha\) is the sheaf-level associator (\cref{cor:sheafTensorObjAssoc})
  and the leftmost arrow is the index-equality isomorphism induced by
  \((n_a+n_b)+n_c = n_a+(n_b+n_c)\) (\cref{def:sheafTensorPow}). This is the
  associativity constraint asserted in \cref{lem:sheafTensorPow_add}.
\end{lemma}

\begin{proof}
  \uses{lem:sheafTensorPow_add, cor:sheafTensorObjAssoc,
    lem:tensorPowAdd_zero, def:modulesLocalizedMonoidal,
    lem:monoidal_pentagon_mathlib, lem:monoidal_triangle_mathlib,
    lem:tensorObjAssoc_eq_localizedAssociator, lem:tensorObjUnitor_eq_localized}
  Induct on \(n_a\). The base case \(n_a = 0\) reduces, via
  \cref{lem:tensorPowAdd_zero}, to the triangle identity relating the left unitor
  to the associator, which is the triangle law
  \cref{lem:monoidal_triangle_mathlib} of the localized structure
  \cref{def:modulesLocalizedMonoidal} (with the unitor bridge
  \cref{lem:tensorObjUnitor_eq_localized}). For the successor step, expand both
  \(\mu_{(k+1)+n_b,n_c}\) and \(\mu_{k+1,\,n_b+n_c}\) by the successor branch of
  \cref{lem:sheafTensorPow_add} and apply the inductive hypothesis to the
  comparison at first index \(k\); the resulting equality of composites is, after
  cancelling index isomorphisms, the image of a pentagon-shaped coherence diagram
  built from associators. That diagram commutes in any monoidal category; the
  localized structure of \cref{def:modulesLocalizedMonoidal} is one on
  \(X.\Modules\), so it commutes there by the pentagon law
  \cref{lem:monoidal_pentagon_mathlib}, the associator bridge
  \cref{lem:tensorObjAssoc_eq_localizedAssociator} rewriting the hand-built
  associator \(\alpha\) (\cref{cor:sheafTensorObjAssoc}) as its localized
  counterpart. This establishes the identity for \(n_a = k+1\) and closes the
  induction.
\end{proof}

\section{Lax-monoidal global sections}
\label{sec:sgr_lax_sections}

The graded multiplication is supplied by the lax-monoidal structure of the
global-sections functor \(\Gamma(X, -) : \ShMod(\mathcal{O}_X) \to
\mathrm{ModuleCat}(\Gamma(X, \mathcal{O}_X))\). It is only \emph{lax} because
global sections do not commute with sheafification: a global section of a
sheafified tensor product is not in general a finite sum of elementary tensors of
global sections.

\begin{definition}\leanok
[Section multiplication]
  \label{def:sectionMul}
  \lean{AlgebraicGeometry.Scheme.Modules.sectionsMul}
  \uses{def:sheafTensorObj, lem:presheafModule_monoidal_mathlib,
    lem:presheafModule_sheafification_mathlib}
  Let \(\mathcal{F}, \mathcal{G}\) be sheaves of \(\mathcal{O}_X\)-modules. The
  \emph{section multiplication} is the \(\Gamma(X,\mathcal{O}_X)\)-bilinear map
  \[
    \Gamma(X, \mathcal{F}) \otimes_{\Gamma(X,\mathcal{O}_X)}
      \Gamma(X, \mathcal{G})
    \;\longrightarrow\;
    \Gamma(X, \mathcal{F} \otimes_{\mathcal{O}_X} \mathcal{G})
  \]
  defined as follows. A pair of global sections \((\sigma, \tau)\) determines,
  by the objectwise formula of \cref{lem:presheafModule_monoidal_mathlib} at the
  top open, the elementary tensor \(\sigma \otimes \tau \in
  (\mathcal{F} \otimes_{p,\mathcal{O}_X} \mathcal{G})(X)\), a global section of
  the tensor product presheaf. Composing with the global-sections component of
  the sheafification unit
  \(\eta : \mathcal{F} \otimes_{p,\mathcal{O}_X} \mathcal{G} \to
  (\mathcal{F} \otimes_{p,\mathcal{O}_X} \mathcal{G})^{\#} =
  \mathcal{F} \otimes_{\mathcal{O}_X} \mathcal{G}\)
  (\cref{lem:presheafModule_sheafification_mathlib},
  \cref{def:sheafTensorObj}) yields a global section
  \(\sigma \cdot \tau \in
  \Gamma(X, \mathcal{F} \otimes_{\mathcal{O}_X} \mathcal{G})\). Bilinearity over
  \(\Gamma(X,\mathcal{O}_X)\) is inherited from bilinearity of the elementary
  tensor and \(\Gamma(X,\mathcal{O}_X)\)-linearity of \(\eta\).
\end{definition}

The graded-monoid axioms below are not raw equalities between sections of two
different tensor powers --- on \(\mathbb{N}\) the indices \((i+j)+k\) and
\(i+(j+k)\), or \(0+n\) and \(n\), are equal but not definitionally so, so the two
sides a priori inhabit distinct section modules. Following the convention of Mathlib's tensor-power construction, we phrase each coherence as an
equality in a single section module, obtained by transporting one side
along an index-equality isomorphism, and provide a bridge that repackages such a
transported equality as an equality of dependent pairs.

\begin{definition}\leanok
[Index-equality transport of section components]
  \label{def:sectionsCast}
  \lean{AlgebraicGeometry.Scheme.Modules.sectionsCast}
  \uses{def:sheafTensorPow}
  Let \(\mathcal{L}\) be an invertible sheaf on \(X\) and let \(h : i = j\) be an
  equality of natural numbers. Applying the global-sections functor
  \(\Gamma(X,-)\) to the canonical isomorphism
  \(\mathcal{L}^{\otimes i} \xrightarrow{\sim} \mathcal{L}^{\otimes j}\) induced by
  \(h\) under the tensor-power functor (\cref{def:sheafTensorPow}) yields the
  \(\Gamma(X,\mathcal{O}_X)\)-linear isomorphism
  \[
    \mathrm{sectionsCast}(h) :
    \Gamma(X, \mathcal{L}^{\otimes i}) \;\xrightarrow{\ \sim\ }\;
    \Gamma(X, \mathcal{L}^{\otimes j}),
  \]
  the \emph{section-component transport} along \(h\). It is the section-level
  analogue of Mathlib's tensor-power index cast, which transports
  \(\bigotimes^{i} M\) to \(\bigotimes^{j} M\) along \(i = j\). Because the
  isomorphism induced by the reflexive equality is the identity and
  \(\Gamma(X,-)\) preserves identities, the transport along a reflexive index
  equality is the identity map (the reflexivity simplification
  \(\mathrm{sectionsCast\_refl}\), \cref{lem:sectionsCast_refl}).
\end{definition}

\begin{lemma}\leanok
[Reflexivity of the section transport]
  \label{lem:sectionsCast_refl}
  \lean{AlgebraicGeometry.Scheme.Modules.sectionsCast_refl}
  \uses{def:sectionsCast}
  For any \(i\), the transport along the reflexive equality \(i = i\)
  (\cref{def:sectionsCast}) equals the identity automorphism of
  \(\Gamma(X, \mathcal{L}^{\otimes i})\). This is the section-level counterpart of
  Mathlib's tensor-power reflexivity lemma, and serves as a simplification lemma
  collapsing the transport along a reflexive index equality.
\end{lemma}

\begin{proof}

\leanok
  The canonical isomorphism
  \(\mathcal{L}^{\otimes i} \cong \mathcal{L}^{\otimes i}\) induced by the
  reflexive equality is the identity, and the global-sections functor carries
  identities to identities.
\end{proof}

\begin{lemma}\leanok
[Cast-mediated equality in the graded monoid]
  \label{lem:gradedMonoid_eq_of_cast}
  \lean{AlgebraicGeometry.Scheme.Modules.gradedMonoid_eq_of_cast}
  \uses{def:sectionsCast, lem:sectionsCast_refl}
  Write the carrier of the graded monoid as the dependent sum
  \(\coprod_{n} \Gamma(X, \mathcal{L}^{\otimes n})\), whose elements are pairs
  \((n, s)\) with \(s \in \Gamma(X, \mathcal{L}^{\otimes n})\). Let
  \(a = (i, x)\) and \(b = (j, y)\) be two such pairs. If \(h : i = j\) and the
  transported component agrees, \(\mathrm{sectionsCast}(h)\,x = y\), then
  \(a = b\). That is, a component-level equality mediated by the index transport
  \cref{def:sectionsCast} repackages into a genuine equality of dependent pairs.
  This is the section-level analogue of Mathlib's bridge from a transport-mediated
  component equality to an equality of dependent pairs in the graded monoid.
\end{lemma}

\begin{proof}
  \uses{lem:sectionsCast_refl}
  \leanok
  Substituting \(j = i\) by means of \(h\) reduces the transport to the identity
  (\cref{lem:sectionsCast_refl}), so the hypothesis becomes \(x = y\); the two
  dependent pairs \((i, x)\) and \((i, y)\) then coincide.
\end{proof}

The four component-level identities of \cref{lem:sectionMul_coherent} are each
proved by reducing the degreewise multiplication to a single section-level
\emph{core} identity, which records how the global-sections functor
\(\Gamma(X,-)\) interacts with one structural isomorphism of the sheaf tensor
product applied to a sheafification-unit image. Two purely bookkeeping helpers
are isolated first, then the four cores.

\begin{lemma}\leanok
[Global sections of a composite]
  \label{lem:val_app_top_comp}
  \lean{AlgebraicGeometry.Scheme.Modules.val_app_top_comp}
  \uses{def:sheafTensorObj}
  For morphisms \(f : \mathcal{A} \to \mathcal{B}\) and
  \(g : \mathcal{B} \to \mathcal{C}\) of sheaves of \(\mathcal{O}_X\)-modules and a
  global section \(x \in \Gamma(X, \mathcal{A})\),
  \[
    \Gamma(g \circ f)(x) \;=\; \Gamma(g)\bigl(\Gamma(f)(x)\bigr),
  \]
  i.e.\ evaluation at the top open of the underlying presheaf morphism turns
  composition into composition.
\end{lemma}

\begin{proof}

\leanok
  Functoriality of the underlying-presheaf evaluation at the top open: the
  value at \(\top\) of \(g \circ f\) is the composite of the values, and applying
  it to \(x\) gives the iterated application.
\end{proof}

\begin{lemma}\leanok
[Decomposition of the left unitor]
  \label{lem:tensorObjUnitIso_hom}
  \lean{AlgebraicGeometry.Scheme.Modules.tensorObjUnitIso_hom}
  \uses{def:tensorObjUnitIso, def:schemeModuleSheafification}
  For a sheaf of \(\mathcal{O}_X\)-modules \(\mathcal{G}\), the forward map of the
  left unitor (\cref{def:tensorObjUnitIso}) is the sheafified presheaf left unitor
  followed by the sheafification counit:
  \[
    \lambda_{\mathcal{G}}
      = \varepsilon_{\mathcal{G}} \circ (\lambda^p_{\mathcal{G}})^{\#},
  \]
  where \(\lambda^p\) is the presheaf left unitor, \((-)^{\#}\) is module
  sheafification (\cref{def:schemeModuleSheafification}) and
  \(\varepsilon_{\mathcal{G}}\) is the reflective counit isomorphism identifying
  the sheafification of the underlying presheaf of \(\mathcal{G}\) with
  \(\mathcal{G}\).
\end{lemma}

\begin{proof}

\leanok
  This is the definition of the left unitor (\cref{def:tensorObjUnitIso}) read on
  forward maps.
\end{proof}

\begin{lemma}

\leanok
[Decomposition of the right unitor]
  \label{lem:tensorObjRightUnitor_hom}
  \lean{AlgebraicGeometry.Scheme.Modules.tensorObjRightUnitor_hom}
  \uses{def:tensorObjRightUnitor, def:tensorBraiding, def:tensorObjUnitIso}
  For a sheaf of \(\mathcal{O}_X\)-modules \(\mathcal{G}\), the forward map of the
  right unitor (\cref{def:tensorObjRightUnitor}) is the braiding followed by the
  left unitor:
  \[
    (\rho_{\mathcal{G}})_{\mathrm{hom}}
      = (\beta_{\mathcal{G},\mathbf{1}_X})_{\mathrm{hom}}
        \,;\, (\lambda_{\mathcal{G}})_{\mathrm{hom}},
  \]
  i.e.\ \(\rho_{\mathcal{G}} = \lambda_{\mathcal{G}} \circ
  \beta_{\mathcal{G},\mathbf{1}_X}\), where \(\beta\) is the braiding
  (\cref{def:tensorBraiding}) and \(\lambda\) the left unitor
  (\cref{def:tensorObjUnitIso}).
\end{lemma}

\begin{proof}

\leanok
  This is the definition of the right unitor (\cref{def:tensorObjRightUnitor}) read
  on forward maps.
\end{proof}

\begin{lemma}\leanok
[Left-unitor core of the section multiplication]
  \label{lem:sectionMul_unitor_core}
  \lean{AlgebraicGeometry.Scheme.Modules.unitor_sectionsMul}
  \uses{def:sectionMul, def:tensorObjUnitIso, lem:tensorObjUnitIso_hom,
    lem:val_app_top_comp, def:schemeModuleSheafification,
    lem:presheafModule_monoidal_mathlib,
    lem:presheafModule_sheafification_mathlib}
  For a sheaf of \(\mathcal{O}_X\)-modules \(\mathcal{G}\) and
  \(a \in \Gamma(X, \mathcal{G})\), applying \(\Gamma(X,-)\) of the left unitor
  (\cref{def:tensorObjUnitIso}) to the section multiplication of \(1\) and \(a\)
  returns \(a\):
  \[
    \Gamma\bigl(\lambda_{\mathcal{G}}\bigr)
      \bigl( 1 \cdot a \bigr) = a,
    \qquad
    1 \cdot a = \mathrm{sectionsMul}(\mathbf{1}_X, \mathcal{G})(1 \otimes a)
      \in \Gamma(X, \mathbf{1}_X \otimes_{\mathcal{O}_X} \mathcal{G}),
  \]
  where \(\mathrm{sectionsMul}\) is the section multiplication
  (\cref{def:sectionMul}) and \(1 \in \Gamma(X,\mathcal{O}_X)\). This is the
  section-level adjunction-transpose core of left unitality.
\end{lemma}

\begin{proof}
  \uses{def:sectionMul, lem:tensorObjUnitIso_hom, lem:val_app_top_comp,
    lem:presheafModule_monoidal_mathlib,
    lem:presheafModule_sheafification_mathlib}
    \leanok
  By \cref{lem:tensorObjUnitIso_hom},
  \(\lambda_{\mathcal{G}} = \varepsilon_{\mathcal{G}} \circ
  (\lambda^p_{\mathcal{G}})^{\#}\); splitting the global-sections application along
  this composite (\cref{lem:val_app_top_comp}) reduces the claim to
  \[
    \Gamma(\varepsilon_{\mathcal{G}})
      \Bigl( \Gamma\bigl((\lambda^p_{\mathcal{G}})^{\#}\bigr)
        \bigl( \eta_{\mathbf{1}_X \otimes_p \mathcal{G}}(1 \otimes a) \bigr) \Bigr)
      = a,
  \]
  the section multiplication being the top-open component of the sheafification
  unit \(\eta\) (\cref{def:sectionMul},
  \cref{lem:presheafModule_sheafification_mathlib}). The composite
  \((\lambda^p_{\mathcal{G}})^{\#}\) followed by the counit
  \(\varepsilon_{\mathcal{G}}\) is, by the right-triangle identity of the
  sheafification adjunction, precisely the presheaf left unitor precomposed with
  the unit: \(\varepsilon_{\mathcal{G}} \circ (\lambda^p_{\mathcal{G}})^{\#} \circ
  \eta = \lambda^p_{\mathcal{G}}\). Hence the left side equals
  \(\Gamma(\lambda^p_{\mathcal{G}})(1 \otimes a)\). Finally the presheaf left
  unitor at the top open is the strict module left unitor
  (\cref{lem:presheafModule_monoidal_mathlib}), which sends \(1 \otimes a\) to
  \(1 \cdot a = a\).
\end{proof}

\begin{lemma}\leanok
[Right-unitor core of the section multiplication]
  \label{lem:sectionMul_rightUnitor_core}
  \lean{AlgebraicGeometry.Scheme.Modules.sectionMul_rightUnitor_core}
  \uses{def:sectionMul, def:tensorObjRightUnitor, def:tensorBraiding,
    lem:sectionMul_unitor_core, lem:val_app_top_comp,
    lem:presheafModule_monoidal_mathlib,
    lem:presheafModule_sheafification_mathlib}
  For a sheaf of \(\mathcal{O}_X\)-modules \(\mathcal{G}\) and
  \(a \in \Gamma(X, \mathcal{G})\), applying \(\Gamma(X,-)\) of the right unitor
  (\cref{def:tensorObjRightUnitor}) to the section multiplication of \(a\) and
  \(1\) returns \(a\):
  \[
    \Gamma\bigl(\rho_{\mathcal{G}}\bigr)
      \bigl( \mathrm{sectionsMul}(\mathcal{G}, \mathbf{1}_X)(a \otimes 1) \bigr)
      = a .
  \]
\end{lemma}

\begin{proof}
  \uses{def:sectionMul, def:tensorObjRightUnitor, def:tensorBraiding,
    lem:sectionMul_unitor_core, lem:val_app_top_comp,
    lem:presheafModule_monoidal_mathlib}
    \leanok
  By \cref{def:tensorObjRightUnitor},
  \(\rho_{\mathcal{G}} = \lambda_{\mathcal{G}} \circ
  \beta_{\mathcal{G},\mathbf{1}_X}\); splitting the application
  (\cref{lem:val_app_top_comp}) gives
  \(\Gamma(\rho_{\mathcal{G}})(\cdots) =
  \Gamma(\lambda_{\mathcal{G}})\bigl(\Gamma(\beta_{\mathcal{G},\mathbf{1}_X})
  (\mathrm{sectionsMul}(\mathcal{G},\mathbf{1}_X)(a \otimes 1))\bigr)\). The
  braiding \(\beta_{\mathcal{G},\mathbf{1}_X}\) is the sheafification of the
  presheaf braiding (\cref{def:tensorBraiding}), and the section multiplication is
  the unit \(\eta\) at the top open; by naturality of \(\eta\) against the
  presheaf braiding \(\beta^p\) (\cref{lem:presheafModule_monoidal_mathlib}), whose
  top-open value sends \(a \otimes 1\) to \(1 \otimes a\),
  \[
    \Gamma(\beta_{\mathcal{G},\mathbf{1}_X})
      \bigl(\eta(a \otimes 1)\bigr)
      = \eta\bigl(\beta^p(a \otimes 1)\bigr) = \eta(1 \otimes a)
      = \mathrm{sectionsMul}(\mathbf{1}_X, \mathcal{G})(1 \otimes a) .
  \]
  Applying \(\Gamma(\lambda_{\mathcal{G}})\) and the left-unitor core
  (\cref{lem:sectionMul_unitor_core}) gives \(a\). Only unit-naturality of the
  presheaf braiding is used beyond the left-unitor core; no further counit or
  triangle input is needed.
\end{proof}

\begin{lemma}\leanok
[Braiding core of the section multiplication]
  \label{lem:sectionMul_braiding_core}
  \lean{AlgebraicGeometry.Scheme.Modules.sectionMul_braiding_core}
  \uses{def:sectionMul, def:tensorBraiding, lem:val_app_top_comp,
    lem:presheafModule_monoidal_mathlib,
    lem:presheafModule_sheafification_mathlib}
  For sheaves of \(\mathcal{O}_X\)-modules \(\mathcal{F}, \mathcal{G}\) and global
  sections \(\sigma \in \Gamma(X,\mathcal{F})\),
  \(\tau \in \Gamma(X,\mathcal{G})\), applying \(\Gamma(X,-)\) of the braiding
  (\cref{def:tensorBraiding}) to the section multiplication swaps the factors:
  \[
    \Gamma\bigl(\beta_{\mathcal{F},\mathcal{G}}\bigr)
      \bigl( \mathrm{sectionsMul}(\mathcal{F},\mathcal{G})(\sigma \otimes \tau)\bigr)
      = \mathrm{sectionsMul}(\mathcal{G},\mathcal{F})(\tau \otimes \sigma) .
  \]
\end{lemma}

\begin{proof}
  \uses{def:sectionMul, def:tensorBraiding, lem:presheafModule_monoidal_mathlib,
    lem:presheafModule_sheafification_mathlib}
    \leanok
  The braiding is the sheafification of the presheaf braiding \(\beta^p\)
  (\cref{def:tensorBraiding}) and the section multiplication is the
  sheafification unit \(\eta\) at the top open (\cref{def:sectionMul}). The
  naturality square of \(\eta\) against the presheaf morphism
  \(\beta^p_{\mathcal{F},\mathcal{G}} : \mathcal{F} \otimes_p \mathcal{G} \to
  \mathcal{G} \otimes_p \mathcal{F}\) reads
  \((\beta^p)^{\#} \circ \eta_{\mathcal{F} \otimes_p \mathcal{G}}
  = \eta_{\mathcal{G} \otimes_p \mathcal{F}} \circ \beta^p\); evaluating at the
  top open on \(\sigma \otimes \tau\), and using that \(\beta^p\) sends
  \(\sigma \otimes \tau\) to \(\tau \otimes \sigma\)
  (\cref{lem:presheafModule_monoidal_mathlib}), gives
  \(\Gamma(\beta_{\mathcal{F},\mathcal{G}})(\eta(\sigma \otimes \tau)) =
  \eta(\tau \otimes \sigma)\), which is the right-hand side. Only unit-naturality
  of the presheaf braiding is used.
\end{proof}

\begin{lemma}\leanok
[Associativity core of the section multiplication]
  \label{lem:sectionMul_assoc_core}
  \lean{AlgebraicGeometry.Scheme.Modules.sectionMul_assoc_core}
  \uses{def:sectionMul, cor:sheafTensorObjAssoc,
    lem:isIso_sheafification_whiskerRight_unit, lem:val_app_top_comp,
    lem:presheafModule_monoidal_mathlib,
    lem:presheafModule_sheafification_mathlib}
  For sheaves of \(\mathcal{O}_X\)-modules \(\mathcal{A}, \mathcal{B},
  \mathcal{C}\) and global sections \(\sigma \in \Gamma(X,\mathcal{A})\),
  \(\tau \in \Gamma(X,\mathcal{B})\), \(\upsilon \in \Gamma(X,\mathcal{C})\),
  applying \(\Gamma(X,-)\) of the associator (\cref{cor:sheafTensorObjAssoc}) to
  the left-bracketed iterated section multiplication yields the right-bracketed
  one:
  \[
    \Gamma\bigl(\alpha_{\mathcal{A},\mathcal{B},\mathcal{C}}\bigr)
      \Bigl( \bigl(\sigma \cdot \tau\bigr) \cdot \upsilon \Bigr)
      = \sigma \cdot \bigl(\tau \cdot \upsilon\bigr),
  \]
  where \(\sigma \cdot \tau = \mathrm{sectionsMul}(\mathcal{A},\mathcal{B})
  (\sigma \otimes \tau)\) and the outer products are the section multiplications
  of the respective tensor products (\cref{def:sectionMul}).
\end{lemma}

\begin{proof}
  \uses{def:sectionMul, cor:sheafTensorObjAssoc,
    lem:tensorObjAssoc_eq_localizedAssociator, def:modulesLocalizedMonoidal,
    thm:localizedMonoidal_mathlib, lem:presheafModule_monoidal_mathlib,
    lem:presheafModule_sheafification_mathlib}
  Both sides are iterated images of elementary tensors under the sheafification
  unit \(\eta\) (\cref{def:sectionMul}): the left side is the \(\eta\)-image of
  \((\sigma \otimes \tau) \otimes \upsilon\) and the right side that of
  \(\sigma \otimes (\tau \otimes \upsilon)\), each through the two-step
  sheafification recorded by \cref{def:sheafTensorObj}. By the associator bridge
  \cref{lem:tensorObjAssoc_eq_localizedAssociator}, the hand-built associator
  \(\alpha_{\mathcal{A},\mathcal{B},\mathcal{C}}\) (\cref{cor:sheafTensorObjAssoc})
  coincides with the localized associator \(\alpha^{\mathrm{loc}}\) of
  \cref{def:modulesLocalizedMonoidal}, whose action on such a tensor is computed
  componentwise by \(\mathrm{associator\_hom\_app}\)
  (\cref{thm:localizedMonoidal_mathlib}): the strong-monoidality comparison
  \(\mu\) carries the iterated \(\eta\)-image to the sheafified presheaf associator
  \((\alpha^p)^{\#}\), whose top-open value reassociates
  \((\sigma \otimes \tau) \otimes \upsilon\) to
  \(\sigma \otimes (\tau \otimes \upsilon)\) by naturality of \(\eta\) against the
  presheaf associator \(\alpha^p\) (\cref{lem:presheafModule_monoidal_mathlib}).
  This gives the claim.
\end{proof}

\begin{lemma}[Coherence of the section multiplication]
  \label{lem:sectionMul_coherent}
  \lean{AlgebraicGeometry.Scheme.Modules.sectionsMul_one_mul,
    AlgebraicGeometry.Scheme.Modules.sectionsMul_mul_one,
    AlgebraicGeometry.Scheme.Modules.sectionsMul_mul_assoc,
    AlgebraicGeometry.Scheme.Modules.sectionsMul_mul_comm}
  \uses{def:sectionMul, def:sectionsCast, lem:moduleUnit_mathlib,
    lem:sheafTensorPow_add}
  Write \(a \cdot b\) for the degreewise section multiplication on tensor powers:
  for \(a \in \Gamma(X, \mathcal{L}^{\otimes i})\) and
  \(b \in \Gamma(X, \mathcal{L}^{\otimes j})\) it is the elementary section
  multiplication (\cref{def:sectionMul}) followed by the tensor-power comparison
  isomorphism \(\mu_{i,j}\) (\cref{lem:sheafTensorPow_add}), landing in
  \(\Gamma(X, \mathcal{L}^{\otimes (i+j)})\); and let \(1\) denote the unit of
  degree \(0\), the image of \(1 \in \Gamma(X,\mathcal{O}_X)\) under the
  identification \(\Gamma(X,\mathcal{O}_X) = \Gamma(X, \mathcal{L}^{\otimes 0})\)
  (\cref{lem:moduleUnit_mathlib}). Mirroring Mathlib's unitality, associativity,
  and commutativity identities for the abstract tensor power, the multiplication
  satisfies the following four component-level identities. Each is
  an honest equality in a single section module, obtained by transporting one side
  along the relevant index equality via \cref{def:sectionsCast}; the subscript on
  \(\mathrm{sectionsCast}\) records that index equality.
  \begin{itemize}
    \item \emph{Left unitality.} For \(a \in \Gamma(X, \mathcal{L}^{\otimes n})\),
      \[
        \mathrm{sectionsCast}_{\,0+n=n}\,(1 \cdot a) = a .
      \]
    \item \emph{Right unitality.} For \(a \in \Gamma(X, \mathcal{L}^{\otimes n})\),
      \[
        \mathrm{sectionsCast}_{\,n+0=n}\,(a \cdot 1) = a .
      \]
    \item \emph{Associativity.} For
      \(a \in \Gamma(X, \mathcal{L}^{\otimes n_a})\),
      \(b \in \Gamma(X, \mathcal{L}^{\otimes n_b})\),
      \(c \in \Gamma(X, \mathcal{L}^{\otimes n_c})\),
      \[
        \mathrm{sectionsCast}_{\,(n_a+n_b)+n_c = n_a+(n_b+n_c)}\,
          ((a \cdot b) \cdot c) = a \cdot (b \cdot c) .
      \]
    \item \emph{Commutativity.} For
      \(a \in \Gamma(X, \mathcal{L}^{\otimes n_a})\) and
      \(b \in \Gamma(X, \mathcal{L}^{\otimes n_b})\),
      \[
        \mathrm{sectionsCast}_{\,n_a+n_b = n_b+n_a}\,(a \cdot b) = b \cdot a .
      \]
  \end{itemize}
  The transport is indispensable: the indices on the two sides of each identity
  are equal but not definitionally so, hence a priori live in distinct section
  modules, and \cref{def:sectionsCast} moves one side into the other's module to
  produce a genuine equality there.
\end{lemma}

\begin{proof}
  \uses{def:sectionMul, def:sectionsCast, lem:sheafTensorPow_add,
    lem:tensorPowAdd_zero, lem:tensorPowAdd_rightUnit, lem:tensorPowAdd_braiding,
    lem:tensorPowAdd_assoc, lem:sectionMul_unitor_core,
    lem:sectionMul_rightUnitor_core, lem:sectionMul_braiding_core,
    lem:sectionMul_assoc_core, lem:val_app_top_comp}
  The degreewise multiplication is, by definition, the section multiplication
  \(\mathrm{sectionsMul}\) (\cref{def:sectionMul}) followed by \(\Gamma(X,-)\) of
  the tensor-power comparison \(\mu\) (\cref{lem:sheafTensorPow_add}). Each
  identity is obtained by first \emph{factoring \(\mu\)} along the appropriate
  branch of \cref{lem:sheafTensorPow_add} and then \emph{cancelling} the resulting
  index-equality isomorphism against the outer \(\mathrm{sectionsCast}\)
  (\cref{def:sectionsCast}), which is by construction the image under
  \(\Gamma(X,-)\) of that same index isomorphism; what remains in each case is one
  of the section-level cores proved above. Splitting the global-sections
  application of a composite uses \cref{lem:val_app_top_comp} throughout.

  \emph{Left unitality.} By \cref{lem:tensorPowAdd_zero}, \(\mu_{0,n}\) is the left
  unitor \(\lambda_{\mathcal{L}^{\otimes n}}\) followed by the index isomorphism
  for \(0+n=n\). Cancelling that isomorphism against
  \(\mathrm{sectionsCast}_{\,0+n=n}\) reduces the left-unitality identity to
  \(\Gamma(\lambda_{\mathcal{L}^{\otimes n}})(1 \cdot a) = a\), which is the
  left-unitor core \cref{lem:sectionMul_unitor_core}.

  \emph{Right unitality.} By \cref{lem:tensorPowAdd_rightUnit}, \(\mu_{n,0}\) is
  the right unitor \(\rho_{\mathcal{L}^{\otimes n}}\) followed by the index
  isomorphism for \(n+0=n\). Cancelling against \(\mathrm{sectionsCast}_{\,n+0=n}\)
  reduces right unitality to
  \(\Gamma(\rho_{\mathcal{L}^{\otimes n}})
  (\mathrm{sectionsMul}(\mathcal{L}^{\otimes n}, \mathbf{1}_X)(a \otimes 1)) = a\),
  the right-unitor core \cref{lem:sectionMul_rightUnitor_core}.

  \emph{Commutativity.} By \cref{lem:tensorPowAdd_braiding}, \(\mu_{n_a,n_b}\) is
  the braiding \(\beta_{\mathcal{L}^{\otimes n_a}, \mathcal{L}^{\otimes n_b}}\)
  followed by \(\mu_{n_b,n_a}\) and the index isomorphism for
  \(n_b+n_a = n_a+n_b\). Cancelling that isomorphism against
  \(\mathrm{sectionsCast}_{\,n_a+n_b = n_b+n_a}\) and applying the braiding core
  \cref{lem:sectionMul_braiding_core}, which carries
  \(\Gamma(\beta)\) of \(\mathrm{sectionsMul}(\cdot)(a \otimes b)\) to
  \(\mathrm{sectionsMul}(\cdot)(b \otimes a)\), identifies the two sides as
  \(b \cdot a\).

  \emph{Associativity.} By \cref{lem:tensorPowAdd_assoc}, the two bracketings of
  the triple comparison agree after the index isomorphism for
  \((n_a+n_b)+n_c = n_a+(n_b+n_c)\); cancelling it against
  \(\mathrm{sectionsCast}_{\,(n_a+n_b)+n_c = n_a+(n_b+n_c)}\) and inserting the
  associativity core \cref{lem:sectionMul_assoc_core}, which intertwines the two
  iterated section multiplications through \(\Gamma\) of the associator
  \(\alpha\), identifies \((a \cdot b) \cdot c\) transported with
  \(a \cdot (b \cdot c)\).
\end{proof}

\section{Graded-ring and graded-module assembly}
\label{sec:sgr_graded_assembly}

The components and multiplications of the previous two sections are assembled
into the graded ring and graded module through Mathlib's
\emph{external-direct-sum} graded API: the relevant inputs are families of
abelian groups indexed by \(\mathbb{N}\) together with multiplication maps
landing in the index-sum component, not submodules of a pre-existing ring.

\begin{definition}\leanok
[Section graded ring]
  \label{def:sectionGradedRing}
  \lean{AlgebraicGeometry.Scheme.Modules.sectionDeg}
  \uses{def:sheafTensorPow, lem:sectionGradedRing_gcommSemiring}
  Let \(L\) be a line bundle (invertible sheaf) on \(X\). The \emph{section
  graded ring}
  \[
    R(X, L) \;=\; \bigoplus_{m \ge 0} \Gamma\bigl(X, L^{\otimes m}\bigr)
  \]
  is the external direct sum of the \(\Gamma(X,\mathcal{O}_X)\)-modules
  \(\mathrm{sectionDeg}\,L\,m = \Gamma(X, L^{\otimes m})\) of global sections of
  the tensor powers (\cref{def:sheafTensorPow}), equipped with the graded
  commutative semiring structure of \cref{lem:sectionGradedRing_gcommSemiring}.
  Its underlying graded datum is the family
  \(m \mapsto \mathrm{sectionDeg}\,L\,m\); the ring operations are supplied by the
  \(\mathrm{GCommSemiring}\) instance on that family.
\end{definition}

\begin{definition}\leanok
[Section graded module]
  \label{def:sectionGradedModule}
  \lean{AlgebraicGeometry.Scheme.Modules.moduleTensorPow}
  \uses{def:sheafModuleTwist, lem:sectionGradedModule_gmodule}
  Let \(L\) be a line bundle and \(\mathcal{F}\) a sheaf of modules on \(X\). The
  \emph{section graded module}
  \[
    M(X, \mathcal{F}, L) \;=\; \bigoplus_{m \ge 0}
      \Gamma\bigl(X, \mathcal{F} \otimes L^{\otimes m}\bigr)
  \]
  is the external direct sum of the global sections of the twists
  \(\mathcal{F} \otimes L^{\otimes m}\) (\cref{def:sheafModuleTwist}); it is a
  graded module over the section graded ring \(R(X,L)\)
  (\cref{def:sectionGradedRing}) via the \(\mathrm{Gmodule}\) structure of
  \cref{lem:sectionGradedModule_gmodule}.
\end{definition}

\begin{lemma}[External-direct-sum graded commutative semiring]
  \label{lem:directSum_gcommSemiring_mathlib}
  \lean{DirectSum.GCommSemiring}
  \mathlibok
  \textit{Provided by Mathlib (\texttt{Mathlib.Algebra.DirectSum.Ring}).}
  Let \(A : \mathbb{N} \to \mathrm{Type}\) be a family of additive commutative
  monoids equipped with a graded multiplication
  \(A_i \times A_j \to A_{i+j}\) and a unit \(1 \in A_0\) satisfying graded
  associativity, two-sided unitality, distributivity, and graded commutativity
  (the data of a \(\mathrm{GCommSemiring}\) on \(A\)). Then the external direct
  sum \(\bigoplus_i A_i\) is a commutative semiring, with multiplication the
  bilinear extension of the graded multiplication and identity the image of
  \(1 \in A_0\).
\end{lemma}

\begin{lemma}[External-direct-sum graded module]
  \label{lem:directSum_gmodule_mathlib}
  \lean{DirectSum.Gmodule}
  \mathlibok
  \textit{Provided by Mathlib (\texttt{Mathlib.Algebra.Module.GradedModule}).}
  Let \(A : \mathbb{N} \to \mathrm{Type}\) carry a graded semiring structure as
  in \cref{lem:directSum_gcommSemiring_mathlib} and let
  \(M : \mathbb{N} \to \mathrm{Type}\) be a family of additive commutative
  monoids with a graded scalar action \(A_i \times M_j \to M_{i+j}\) satisfying
  the graded module axioms (the data of a \(\mathrm{Gmodule}\) of \(M\) over
  \(A\)). Then the external direct sum \(\bigoplus_j M_j\) is a module over the
  semiring \(\bigoplus_i A_i\), with scalar multiplication the bilinear
  extension of the graded action.
\end{lemma}

\begin{lemma}\leanok
[Graded semiring structure on the section components]
  \label{lem:sectionGradedRing_gcommSemiring}
  \lean{AlgebraicGeometry.sectionGradedRing_gcommSemiring}
  \uses{def:sheafTensorPow, lem:sheafTensorPow_add, def:sectionMul,
    lem:sectionMul_coherent, def:sectionsCast, lem:gradedMonoid_eq_of_cast,
    lem:directSum_gcommSemiring_mathlib, lem:moduleUnit_mathlib}
  % SOURCE: [Stacks Project], "Sheaves of Modules", \S "Invertible modules",
  % preceding `definition-gamma-star' (Tag 01CV) (read from
  % references/stacks-modules.tex, L4257-L4267).
  % SOURCE QUOTE:
  % "Let $(X, \mathcal{O}_X)$ be a ringed space. We can define a $\mathbf{Z}$-graded
  %  ring structure on $\bigoplus \Gamma(X, \mathcal{L}^{\otimes n})$ by mapping
  %  $s \in \Gamma(X, \mathcal{L}^{\otimes n})$ and
  %  $t \in \Gamma(X, \mathcal{L}^{\otimes m})$ to the
  %  section corresponding to $s \otimes t$ in
  %  $\Gamma(X, \mathcal{L}^{\otimes n + m})$.
  %  We omit the verification that this defines a commutative
  %  and associative ring with $1$."
  \textit{Source: [Stacks Project], "Sheaves of Modules", Tag 01CV (the graded
  ring structure on \(\bigoplus \Gamma(X,\mathcal{L}^{\otimes n})\)).}
  Let \(L_s\) be a line bundle on \(X_s\). The family of
  \(\Gamma(X_s,\mathcal{O}_{X_s})\)-modules
  \(m \mapsto \Gamma(X_s, L_s^{\otimes m})\) carries a graded commutative
  semiring structure (a \(\mathrm{GCommSemiring}\),
  \cref{lem:directSum_gcommSemiring_mathlib}) whose degree-\((m,m')\)
  multiplication is the composite
  \[
    \Gamma(X_s, L_s^{\otimes m}) \otimes_{\Gamma(X_s,\mathcal{O}_{X_s})}
      \Gamma(X_s, L_s^{\otimes m'})
    \xrightarrow{\ \cdot\ }
    \Gamma\bigl(X_s, L_s^{\otimes m} \otimes L_s^{\otimes m'}\bigr)
    \xrightarrow{\ \Gamma(\mu_{m,m'})\ }
    \Gamma\bigl(X_s, L_s^{\otimes (m+m')}\bigr),
  \]
  the section multiplication (\cref{def:sectionMul}) followed by the
  tensor-power comparison isomorphism (\cref{lem:sheafTensorPow_add}), with unit
  \(1 \in \Gamma(X_s, \mathcal{O}_{X_s}) = \Gamma(X_s, L_s^{\otimes 0})\)
  (\cref{lem:moduleUnit_mathlib}). Consequently
  \(R(X_s, L_s) = \bigoplus_{m \ge 0} \Gamma(X_s, L_s^{\otimes m})\) is a
  commutative semiring; this is the graded-ring structure underlying
  \cref{def:sectionGradedRing}.
\end{lemma}

\begin{proof}
  \uses{lem:sheafTensorPow_add, lem:sectionMul_coherent, def:sectionsCast,
    lem:gradedMonoid_eq_of_cast, def:sectionMul,
    lem:directSum_gcommSemiring_mathlib}
  % NOTE: \leanok is absent from this proof because the graded-semiring instance
  % depends transitively on the three coherences of \cref{lem:sectionMul_coherent},
  % which are not yet closed.
  The structure is assembled field-for-field following Mathlib's construction of
  the graded semiring on abstract tensor powers. The
  degreewise multiplication \(\Gamma(\mu_{m,m'}) \circ (\,\cdot\,)\) of the
  statement and the degree-\(0\) unit \(1 \in \Gamma(X, \mathcal{L}^{\otimes 0})\)
  provide the graded multiplication and unit. Their graded-monoid axioms --- left
  and right unitality and associativity --- are obtained from the corresponding
  component-level identities of \cref{lem:sectionMul_coherent} by passing each
  through the bridge \cref{lem:gradedMonoid_eq_of_cast}, which converts a
  transport-mediated equality (\cref{def:sectionsCast}) into the required equality
  of dependent pairs; the graded power operation is the default iterated product
  and contributes no further data. This yields the graded-monoid layer.

  The remaining semiring axioms are the bilinearity (left and right distributivity
  and annihilation by \(0\)) of the graded multiplication and the behaviour of the
  natural-number coercion. Bilinearity is immediate: the degreewise multiplication
  \(\Gamma(\mu_{m,m'}) \circ (\,\cdot\,)\) is \(\Gamma(X,\mathcal{O}_X)\)-bilinear,
  being the composite of the bilinear section multiplication (\cref{def:sectionMul})
  with the linear comparison \(\Gamma(\mu_{m,m'})\)
  (\cref{lem:sheafTensorPow_add}), so each side vanishes on \(0\) and distributes
  over sums. The natural-number coercion sends \(n\) to \(n \cdot 1\), the
  \(n\)-fold sum of the degree-\(0\) unit. This completes the graded-semiring
  layer.

  Finally graded commutativity is the commutativity identity of
  \cref{lem:sectionMul_coherent}, again passed through
  \cref{lem:gradedMonoid_eq_of_cast}. These are precisely the data of a
  \(\mathrm{GCommSemiring}\) on \(m \mapsto \Gamma(X, \mathcal{L}^{\otimes m})\),
  so \cref{lem:directSum_gcommSemiring_mathlib} endows
  \(\bigoplus_{m \ge 0} \Gamma(X_s, L_s^{\otimes m})\) with its commutative-semiring
  structure.
\end{proof}

\begin{lemma}[Graded module structure on the twisted-section components]
  \label{lem:sectionGradedModule_gmodule}
  \lean{AlgebraicGeometry.sectionGradedModule_gmodule}
  \uses{def:sheafModuleTwist, lem:sheafTensorPow_add, def:sectionMul,
    lem:sectionMul_coherent, def:sectionsCast, lem:gradedMonoid_eq_of_cast,
    lem:directSum_gmodule_mathlib, lem:sectionGradedRing_gcommSemiring}
  % SOURCE: [Stacks Project], "Sheaves of Modules", \S "Invertible modules",
  % `definition-gamma-star' (Tag 01CV) (read from
  % references/stacks-modules.tex, L4279-L4286).
  % SOURCE QUOTE:
  % "Given a sheaf of $\mathcal{O}_X$-modules $\mathcal{F}$ we set
  %  $$
  %  \Gamma_*(X, \mathcal{L}, \mathcal{F})
  %  =
  %  \bigoplus\nolimits_{n \in \mathbf{Z}} \Gamma(X,
  %  \mathcal{F} \otimes_{\mathcal{O}_X} \mathcal{L}^{\otimes n})
  %  $$
  %  which we think of as a graded $\Gamma_*(X, \mathcal{L})$-module."
  \textit{Source: [Stacks Project], "Sheaves of Modules", Tag 01CV (the graded
  module \(\Gamma_*(X,\mathcal{L},\mathcal{F})\)).}
  Let \(\mathcal{F}_s\) be a coherent sheaf and \(L_s\) a line bundle on
  \(X_s\). The family
  \(m \mapsto \Gamma(X_s, \mathcal{F}_s \otimes L_s^{\otimes m})\)
  (\cref{def:sheafModuleTwist}) carries a graded module structure (a
  \(\mathrm{Gmodule}\), \cref{lem:directSum_gmodule_mathlib}) over the graded
  semiring of \cref{lem:sectionGradedRing_gcommSemiring}, with degree-\((m,m')\)
  action
  \[
    \Gamma(X_s, L_s^{\otimes m}) \otimes_{\Gamma(X_s,\mathcal{O}_{X_s})}
      \Gamma(X_s, \mathcal{F}_s \otimes L_s^{\otimes m'})
    \xrightarrow{\ \cdot\ }
    \Gamma\bigl(X_s, L_s^{\otimes m} \otimes
      (\mathcal{F}_s \otimes L_s^{\otimes m'})\bigr)
    \xrightarrow{\ \sim\ }
    \Gamma\bigl(X_s, \mathcal{F}_s \otimes L_s^{\otimes (m+m')}\bigr),
  \]
  the section multiplication followed by the comparison isomorphism that
  reassociates and merges the line-bundle factors (using
  \cref{lem:sheafTensorPow_add} and the symmetry of the tensor product).
  Consequently \(M(X_s, \mathcal{F}_s, L_s) = \bigoplus_{m \ge 0}
  \Gamma(X_s, \mathcal{F}_s \otimes L_s^{\otimes m})\) is a graded module over
  \(R(X_s, L_s)\); this is the graded-module structure underlying
  \cref{def:sectionGradedModule}.
\end{lemma}

\begin{proof}
  \uses{lem:sectionMul_coherent, lem:sheafTensorPow_add, def:sectionsCast,
    lem:gradedMonoid_eq_of_cast, def:sectionMul,
    lem:directSum_gmodule_mathlib, lem:sectionGradedRing_gcommSemiring}
  The action maps are the displayed composites: the section multiplication
  (\cref{def:sectionMul}) followed by the comparison isomorphism that reassociates
  and merges the line-bundle factors (\cref{lem:sheafTensorPow_add}). The
  \(\mathrm{Gmodule}\) is assembled field-for-field as the graded-action precedent
  \(\mathrm{GradedMonoid.GMulAction}\) of Mathlib, with the action grading on
  indices given by ordinary addition on \(\mathbb{N}\) (so the degree-\((m,m')\)
  action lands in degree \(m+m'\)). Its sigma-type axioms --- unitality
  \(1 \cdot x = x\) and compatibility
  \(r \cdot (r' \cdot x) = (r r') \cdot x\) with the ring multiplication --- are
  the unit and associativity coherences of the section multiplication
  (\cref{lem:sectionMul_coherent}) read with the third tensor factor carrying
  \(\mathcal{F}_s\), each converted from its transport-mediated form
  (\cref{def:sectionsCast}) into the required equality of dependent pairs by the
  bridge \cref{lem:gradedMonoid_eq_of_cast}. The bilinearity axioms of the action
  (additivity and annihilation by \(0\) in each variable) are free, the action map
  being \(\Gamma(X,\mathcal{O}_X)\)-bilinear exactly as in the proof of
  \cref{lem:sectionGradedRing_gcommSemiring}. These are the data of a
  \(\mathrm{Gmodule}\), so \cref{lem:directSum_gmodule_mathlib} produces the module
  structure of \(M(X_s, \mathcal{F}_s, L_s)\) over \(R(X_s, L_s)\).
\end{proof}
